% Options for packages loaded elsewhere
\PassOptionsToPackage{unicode}{hyperref}
\PassOptionsToPackage{hyphens}{url}
%
\documentclass[
  11pt,
  a4paper,
]{article}
\usepackage{amsmath,amssymb}
\usepackage{iftex}
\ifPDFTeX
  \usepackage[T1]{fontenc}
  \usepackage[utf8]{inputenc}
  \usepackage{textcomp} % provide euro and other symbols
\else % if luatex or xetex
  \usepackage{unicode-math} % this also loads fontspec
  \defaultfontfeatures{Scale=MatchLowercase}
  \defaultfontfeatures[\rmfamily]{Ligatures=TeX,Scale=1}
\fi
\usepackage{lmodern}
\ifPDFTeX\else
  % xetex/luatex font selection
  \setmainfont[]{Tinos}
  \setmonofont[]{DejaVu Sans Mono}
\fi
% Use upquote if available, for straight quotes in verbatim environments
\IfFileExists{upquote.sty}{\usepackage{upquote}}{}
\IfFileExists{microtype.sty}{% use microtype if available
  \usepackage[]{microtype}
  \UseMicrotypeSet[protrusion]{basicmath} % disable protrusion for tt fonts
}{}
\makeatletter
\@ifundefined{KOMAClassName}{% if non-KOMA class
  \IfFileExists{parskip.sty}{%
    \usepackage{parskip}
  }{% else
    \setlength{\parindent}{0pt}
    \setlength{\parskip}{6pt plus 2pt minus 1pt}}
}{% if KOMA class
  \KOMAoptions{parskip=half}}
\makeatother
\usepackage{xcolor}
\usepackage[top=2cm,bottom=2cm,left=2cm,right=2cm]{geometry}
\usepackage{color}
\usepackage{fancyvrb}
\newcommand{\VerbBar}{|}
\newcommand{\VERB}{\Verb[commandchars=\\\{\}]}
\DefineVerbatimEnvironment{Highlighting}{Verbatim}{commandchars=\\\{\}}
% Add ',fontsize=\small' for more characters per line
\newenvironment{Shaded}{}{}
\newcommand{\AlertTok}[1]{\textcolor[rgb]{1.00,0.00,0.00}{\textbf{#1}}}
\newcommand{\AnnotationTok}[1]{\textcolor[rgb]{0.38,0.63,0.69}{\textbf{\textit{#1}}}}
\newcommand{\AttributeTok}[1]{\textcolor[rgb]{0.49,0.56,0.16}{#1}}
\newcommand{\BaseNTok}[1]{\textcolor[rgb]{0.25,0.63,0.44}{#1}}
\newcommand{\BuiltInTok}[1]{\textcolor[rgb]{0.00,0.50,0.00}{#1}}
\newcommand{\CharTok}[1]{\textcolor[rgb]{0.25,0.44,0.63}{#1}}
\newcommand{\CommentTok}[1]{\textcolor[rgb]{0.38,0.63,0.69}{\textit{#1}}}
\newcommand{\CommentVarTok}[1]{\textcolor[rgb]{0.38,0.63,0.69}{\textbf{\textit{#1}}}}
\newcommand{\ConstantTok}[1]{\textcolor[rgb]{0.53,0.00,0.00}{#1}}
\newcommand{\ControlFlowTok}[1]{\textcolor[rgb]{0.00,0.44,0.13}{\textbf{#1}}}
\newcommand{\DataTypeTok}[1]{\textcolor[rgb]{0.56,0.13,0.00}{#1}}
\newcommand{\DecValTok}[1]{\textcolor[rgb]{0.25,0.63,0.44}{#1}}
\newcommand{\DocumentationTok}[1]{\textcolor[rgb]{0.73,0.13,0.13}{\textit{#1}}}
\newcommand{\ErrorTok}[1]{\textcolor[rgb]{1.00,0.00,0.00}{\textbf{#1}}}
\newcommand{\ExtensionTok}[1]{#1}
\newcommand{\FloatTok}[1]{\textcolor[rgb]{0.25,0.63,0.44}{#1}}
\newcommand{\FunctionTok}[1]{\textcolor[rgb]{0.02,0.16,0.49}{#1}}
\newcommand{\ImportTok}[1]{\textcolor[rgb]{0.00,0.50,0.00}{\textbf{#1}}}
\newcommand{\InformationTok}[1]{\textcolor[rgb]{0.38,0.63,0.69}{\textbf{\textit{#1}}}}
\newcommand{\KeywordTok}[1]{\textcolor[rgb]{0.00,0.44,0.13}{\textbf{#1}}}
\newcommand{\NormalTok}[1]{#1}
\newcommand{\OperatorTok}[1]{\textcolor[rgb]{0.40,0.40,0.40}{#1}}
\newcommand{\OtherTok}[1]{\textcolor[rgb]{0.00,0.44,0.13}{#1}}
\newcommand{\PreprocessorTok}[1]{\textcolor[rgb]{0.74,0.48,0.00}{#1}}
\newcommand{\RegionMarkerTok}[1]{#1}
\newcommand{\SpecialCharTok}[1]{\textcolor[rgb]{0.25,0.44,0.63}{#1}}
\newcommand{\SpecialStringTok}[1]{\textcolor[rgb]{0.73,0.40,0.53}{#1}}
\newcommand{\StringTok}[1]{\textcolor[rgb]{0.25,0.44,0.63}{#1}}
\newcommand{\VariableTok}[1]{\textcolor[rgb]{0.10,0.09,0.49}{#1}}
\newcommand{\VerbatimStringTok}[1]{\textcolor[rgb]{0.25,0.44,0.63}{#1}}
\newcommand{\WarningTok}[1]{\textcolor[rgb]{0.38,0.63,0.69}{\textbf{\textit{#1}}}}
\usepackage{longtable,booktabs,array}
\usepackage{calc} % for calculating minipage widths
% Correct order of tables after \paragraph or \subparagraph
\usepackage{etoolbox}
\makeatletter
\patchcmd\longtable{\par}{\if@noskipsec\mbox{}\fi\par}{}{}
\makeatother
% Allow footnotes in longtable head/foot
\IfFileExists{footnotehyper.sty}{\usepackage{footnotehyper}}{\usepackage{footnote}}
\makesavenoteenv{longtable}
\setlength{\emergencystretch}{3em} % prevent overfull lines
\providecommand{\tightlist}{%
  \setlength{\itemsep}{0pt}\setlength{\parskip}{0pt}}
\setcounter{secnumdepth}{-\maxdimen} % remove section numbering
\usepackage{fancyhdr}
\usepackage{etoolbox}
\AtBeginEnvironment{longtable}{\setlength{\parskip}{0pt}}
\pagestyle{fancy}
\fancyhf{}
\setlength{\headheight}{16pt}
\lhead{\footnotesize\textbf{DSAI1003 -- Fundamental Programming Concepts in Python}}
\rhead{\footnotesize\textbf{Lecture 3: Decisions}}
\lfoot{\scriptsize Faculty of Data Science and Artificial Intelligence -- National Economics University}
\rfoot{\footnotesize Page \thepage}
\renewcommand{\headrulewidth}{0.4pt}
\renewcommand{\footrulewidth}{0.3pt}
\setlength{\emergencystretch}{3em}
\ifLuaTeX
  \usepackage{selnolig}  % disable illegal ligatures
\fi
\usepackage{bookmark}
\IfFileExists{xurl.sty}{\usepackage{xurl}}{} % add URL line breaks if available
\urlstyle{same}
\hypersetup{
  hidelinks,
  pdfcreator={LaTeX via pandoc}}

\title{Lesson 3: Decisions}
\author{Faculty of Data Science \& Artificial Intelligence (FDA) - National Economics University}
\date{Course Code: DSAI1003}

\begin{document}
\maketitle

{
\setcounter{tocdepth}{3}
\tableofcontents
}
\section{Lesson 3. Decisions}\label{lesson-3.-decisions}

\textbf{Last updated:} September 19, 2026

\begin{quote}
\textbf{Source:} Chapter 3 -- \emph{Decisions}, in the uploaded
\emph{Python for Everyone} textbook by Cay Horstmann and Rance Necaise
(2/e).

This lecture follows the structure, terminology, examples, and
instructional focus of Chapter 3 while reorganizing the material for
direct classroom use. The examples use the textbook's naming conventions
and, where formatting is needed, its \texttt{\%}-style string
formatting.
\end{quote}

\begin{center}\rule{0.5\linewidth}{0.5pt}\end{center}

\subsection{Lesson Introduction}\label{lesson-introduction}

The programs in Lesson 2 executed statements mostly from top to bottom.
That is enough for straightforward calculations, but useful programs
must also be able to \textbf{make decisions}.

A decision allows a program to choose different actions depending on
data or circumstances. For example, a program may need to:

\begin{itemize}
\tightlist
\item
  apply one discount for small purchases and another discount for large
  purchases;
\item
  reject an invalid floor number in an elevator program;
\item
  classify an earthquake by magnitude;
\item
  determine whether a string has the required form;
\item
  choose a tax calculation according to marital status and income;
\item
  check whether several conditions are simultaneously true.
\end{itemize}

The central idea of this lesson is:

\begin{Shaded}
\begin{Highlighting}[]
\NormalTok{Evaluate a condition}
\NormalTok{        ↓}
\NormalTok{   True or False?}
\NormalTok{      ↙      ↘}
\NormalTok{ Action A   Action B}
\NormalTok{      ↘      ↙}
\NormalTok{ Continue the program}
\end{Highlighting}
\end{Shaded}

The chapter develops this idea gradually, beginning with a two-way
\texttt{if} statement and then extending it to nested decisions,
multiple alternatives, Boolean expressions, string analysis, systematic
testing, and input validation.

\begin{center}\rule{0.5\linewidth}{0.5pt}\end{center}

\subsection{Knowledge and Skills to Be
Achieved}\label{knowledge-and-skills-to-be-achieved}

After completing this lesson, students will be able to:

\begin{itemize}
\tightlist
\item
  Use \texttt{if}, \texttt{if/else}, and \texttt{if/elif/else}
  statements.
\item
  Explain the role of a \textbf{condition} in program control flow.
\item
  Use indentation correctly to define statement blocks.
\item
  Compare numeric and string values using relational operators.
\item
  Distinguish assignment \texttt{=} from equality testing \texttt{==}.
\item
  Avoid exact equality tests for computed floating-point values when
  approximation is required.
\item
  Design a two-way decision systematically from a problem statement.
\item
  Implement \textbf{nested branches} for multi-level decisions.
\item
  Use \texttt{elif} for mutually exclusive alternatives.
\item
  Order conditions correctly from more specific to more general cases
  when necessary.
\item
  Draw and interpret simple \textbf{flowcharts}.
\item
  Design \textbf{test cases} that cover branches and boundary values.
\item
  Use Boolean values \texttt{True} and \texttt{False}.
\item
  Combine conditions using \texttt{and}, \texttt{or}, and \texttt{not}.
\item
  Explain operator precedence and short-circuit evaluation.
\item
  Apply De Morgan's laws to simplify negated Boolean expressions.
\item
  Analyze strings using membership operators and string methods.
\item
  Validate user input before processing it.
\item
  Normalize text input with \texttt{upper()} or \texttt{lower()} when
  case should not matter.
\end{itemize}

\begin{center}\rule{0.5\linewidth}{0.5pt}\end{center}

\subsection{Lesson Structure}\label{lesson-structure}

\begin{enumerate}
\def\labelenumi{\arabic{enumi}.}
\tightlist
\item
  The \texttt{if} Statement
\item
  Relational Operators
\item
  Nested Branches
\item
  Multiple Alternatives
\item
  Problem Solving: Flowcharts
\item
  Problem Solving: Test Cases
\item
  Boolean Variables and Operators
\item
  Analyzing Strings
\item
  Application: Input Validation
\item
  Optional Topics and Toolboxes
\item
  Summary and Exercises
\end{enumerate}

\begin{center}\rule{0.5\linewidth}{0.5pt}\end{center}

\section{\texorpdfstring{3.1 The \texttt{if}
Statement}{3.1 The if Statement}}\label{the-if-statement}

Programs often need to perform one action when a condition is true and
another action when it is false.

The \texttt{if} statement provides this branching behavior.

\subsection{3.1.1 A Two-Way Decision}\label{a-two-way-decision}

Consider an elevator system in a building that labels the floor after 12
as floor 14. Internally, the physical floor numbers still continue
normally, so a request for displayed floor 20 corresponds to actual
floor 19.

A decision can be written as:

\begin{Shaded}
\begin{Highlighting}[]
\ControlFlowTok{if}\NormalTok{ floor }\OperatorTok{\textgreater{}} \DecValTok{13}\NormalTok{ :}
\NormalTok{    actualFloor }\OperatorTok{=}\NormalTok{ floor }\OperatorTok{{-}} \DecValTok{1}
\ControlFlowTok{else}\NormalTok{ :}
\NormalTok{    actualFloor }\OperatorTok{=}\NormalTok{ floor}
\end{Highlighting}
\end{Shaded}

Conceptually:

\begin{Shaded}
\begin{Highlighting}[]
\NormalTok{                 floor \textgreater{} 13 ?}
\NormalTok{                  /       \textbackslash{}}
\NormalTok{              True       False}
\NormalTok{               /            \textbackslash{}}
\NormalTok{ actualFloor = floor {-} 1   actualFloor = floor}
\end{Highlighting}
\end{Shaded}

Only \textbf{one branch} executes.

If \texttt{floor\ \textgreater{}\ 13} is true, Python executes the first
block and skips the \texttt{else} block. Otherwise, Python skips the
first block and executes the \texttt{else} block.

\begin{center}\rule{0.5\linewidth}{0.5pt}\end{center}

\subsection{3.1.2 General Syntax}\label{general-syntax}

\begin{Shaded}
\begin{Highlighting}[]
\ControlFlowTok{if}\NormalTok{ condition :}
\NormalTok{    statements1}
\ControlFlowTok{else}\NormalTok{ :}
\NormalTok{    statements2}
\end{Highlighting}
\end{Shaded}

The \texttt{else} branch is optional:

\begin{Shaded}
\begin{Highlighting}[]
\ControlFlowTok{if}\NormalTok{ condition :}
\NormalTok{    statements}
\end{Highlighting}
\end{Shaded}

Example:

\begin{Shaded}
\begin{Highlighting}[]
\NormalTok{actualFloor }\OperatorTok{=}\NormalTok{ floor}

\ControlFlowTok{if}\NormalTok{ floor }\OperatorTok{\textgreater{}} \DecValTok{13}\NormalTok{ :}
\NormalTok{    actualFloor }\OperatorTok{=}\NormalTok{ actualFloor }\OperatorTok{{-}} \DecValTok{1}
\end{Highlighting}
\end{Shaded}

This version is useful when no special action is required for the false
case.

\begin{center}\rule{0.5\linewidth}{0.5pt}\end{center}

\subsection{3.1.3 Compound Statements and
Blocks}\label{compound-statements-and-blocks}

An \texttt{if} statement is a \textbf{compound statement}. Its header
ends with a colon, and the statements belonging to the branch are
indented.

\begin{Shaded}
\begin{Highlighting}[]
\ControlFlowTok{if}\NormalTok{ totalSales }\OperatorTok{\textgreater{}} \FloatTok{100.0}\NormalTok{ :}
\NormalTok{    discount }\OperatorTok{=}\NormalTok{ totalSales }\OperatorTok{*} \FloatTok{0.05}
\NormalTok{    totalSales }\OperatorTok{=}\NormalTok{ totalSales }\OperatorTok{{-}}\NormalTok{ discount}
    \BuiltInTok{print}\NormalTok{(}\StringTok{"You received a discount of"}\NormalTok{, discount)}
\end{Highlighting}
\end{Shaded}

The three indented statements form one block.

Important rules:

\begin{itemize}
\tightlist
\item
  The header ends with \texttt{:}.
\item
  Statements in the same block must be indented to the same level.
\item
  The \texttt{if} and its matching \texttt{else} must be aligned.
\item
  The block ends when the indentation returns to an outer level.
\end{itemize}

Python uses indentation as part of its syntax, not merely as decoration.

\begin{center}\rule{0.5\linewidth}{0.5pt}\end{center}

\subsection{Example -- Elevator
Simulation}\label{example-elevator-simulation}

\begin{Shaded}
\begin{Highlighting}[]
\CommentTok{\#\#}
\CommentTok{\# This program simulates an elevator panel that skips floor 13.}
\CommentTok{\#}

\NormalTok{floor }\OperatorTok{=} \BuiltInTok{int}\NormalTok{(}\BuiltInTok{input}\NormalTok{(}\StringTok{"Floor: "}\NormalTok{))}

\ControlFlowTok{if}\NormalTok{ floor }\OperatorTok{\textgreater{}} \DecValTok{13}\NormalTok{ :}
\NormalTok{    actualFloor }\OperatorTok{=}\NormalTok{ floor }\OperatorTok{{-}} \DecValTok{1}
\ControlFlowTok{else}\NormalTok{ :}
\NormalTok{    actualFloor }\OperatorTok{=}\NormalTok{ floor}

\BuiltInTok{print}\NormalTok{(}\StringTok{"The elevator will travel to the actual floor"}\NormalTok{, actualFloor)}
\end{Highlighting}
\end{Shaded}

Example run:

\begin{Shaded}
\begin{Highlighting}[]
\NormalTok{Floor: 20}
\NormalTok{The elevator will travel to the actual floor 19}
\end{Highlighting}
\end{Shaded}

\begin{center}\rule{0.5\linewidth}{0.5pt}\end{center}

\section{Common Error 3.1 -- Tabs and Inconsistent
Indentation}\label{common-error-3.1-tabs-and-inconsistent-indentation}

Python requires statements in one block to have consistent indentation.

Problematic code can look aligned in an editor but internally mix tab
characters and spaces.

Prefer configuring the editor to insert spaces when the Tab key is
pressed.

A common convention is four spaces per indentation level:

\begin{Shaded}
\begin{Highlighting}[]
\ControlFlowTok{if}\NormalTok{ score }\OperatorTok{\textgreater{}=} \DecValTok{50}\NormalTok{ :}
    \BuiltInTok{print}\NormalTok{(}\StringTok{"Pass"}\NormalTok{)}
    \BuiltInTok{print}\NormalTok{(}\StringTok{"Continue to the next course"}\NormalTok{)}
\end{Highlighting}
\end{Shaded}

The important point is consistency.

\begin{center}\rule{0.5\linewidth}{0.5pt}\end{center}

\section{Programming Tip 3.1 -- Avoid Duplication in
Branches}\label{programming-tip-3.1-avoid-duplication-in-branches}

Consider:

\begin{Shaded}
\begin{Highlighting}[]
\ControlFlowTok{if}\NormalTok{ floor }\OperatorTok{\textgreater{}} \DecValTok{13}\NormalTok{ :}
\NormalTok{    actualFloor }\OperatorTok{=}\NormalTok{ floor }\OperatorTok{{-}} \DecValTok{1}
    \BuiltInTok{print}\NormalTok{(}\StringTok{"Actual floor:"}\NormalTok{, actualFloor)}
\ControlFlowTok{else}\NormalTok{ :}
\NormalTok{    actualFloor }\OperatorTok{=}\NormalTok{ floor}
    \BuiltInTok{print}\NormalTok{(}\StringTok{"Actual floor:"}\NormalTok{, actualFloor)}
\end{Highlighting}
\end{Shaded}

The output statement is duplicated.

A clearer version is:

\begin{Shaded}
\begin{Highlighting}[]
\ControlFlowTok{if}\NormalTok{ floor }\OperatorTok{\textgreater{}} \DecValTok{13}\NormalTok{ :}
\NormalTok{    actualFloor }\OperatorTok{=}\NormalTok{ floor }\OperatorTok{{-}} \DecValTok{1}
\ControlFlowTok{else}\NormalTok{ :}
\NormalTok{    actualFloor }\OperatorTok{=}\NormalTok{ floor}

\BuiltInTok{print}\NormalTok{(}\StringTok{"Actual floor:"}\NormalTok{, actualFloor)}
\end{Highlighting}
\end{Shaded}

Moving common work outside the branches:

\begin{itemize}
\tightlist
\item
  reduces duplication;
\item
  makes the program shorter;
\item
  reduces the risk that one copy is changed while another is forgotten.
\end{itemize}

\begin{center}\rule{0.5\linewidth}{0.5pt}\end{center}

\section{Special Topic 3.1 -- Conditional
Expressions}\label{special-topic-3.1-conditional-expressions}

Python can represent a simple two-way choice as an expression:

\begin{Shaded}
\begin{Highlighting}[]
\NormalTok{actualFloor }\OperatorTok{=}\NormalTok{ floor }\OperatorTok{{-}} \DecValTok{1} \ControlFlowTok{if}\NormalTok{ floor }\OperatorTok{\textgreater{}} \DecValTok{13} \ControlFlowTok{else}\NormalTok{ floor}
\end{Highlighting}
\end{Shaded}

This is equivalent to:

\begin{Shaded}
\begin{Highlighting}[]
\ControlFlowTok{if}\NormalTok{ floor }\OperatorTok{\textgreater{}} \DecValTok{13}\NormalTok{ :}
\NormalTok{    actualFloor }\OperatorTok{=}\NormalTok{ floor }\OperatorTok{{-}} \DecValTok{1}
\ControlFlowTok{else}\NormalTok{ :}
\NormalTok{    actualFloor }\OperatorTok{=}\NormalTok{ floor}
\end{Highlighting}
\end{Shaded}

For beginning programs, the full \texttt{if/else} form is often easier
to read, especially when each branch contains more than one action.

\begin{center}\rule{0.5\linewidth}{0.5pt}\end{center}

\subsection{Self Check 3.1}\label{self-check-3.1}

\subsubsection{Question 1}\label{question-1}

What value is stored in \texttt{actualFloor} when \texttt{floor} is
\texttt{18}?

\begin{Shaded}
\begin{Highlighting}[]
\ControlFlowTok{if}\NormalTok{ floor }\OperatorTok{\textgreater{}} \DecValTok{13}\NormalTok{ :}
\NormalTok{    actualFloor }\OperatorTok{=}\NormalTok{ floor }\OperatorTok{{-}} \DecValTok{1}
\ControlFlowTok{else}\NormalTok{ :}
\NormalTok{    actualFloor }\OperatorTok{=}\NormalTok{ floor}
\end{Highlighting}
\end{Shaded}

\textbf{Answer}

\texttt{17}.

\subsubsection{Question 2}\label{question-2}

What value is stored in \texttt{actualFloor} when \texttt{floor} is
\texttt{9}?

\textbf{Answer}

\texttt{9}.

\subsubsection{Question 3}\label{question-3}

Why is the colon required after the \texttt{if} condition?

\textbf{Answer}

It marks the header of a compound statement. The following indented
block belongs to that header.

\subsubsection{Question 4}\label{question-4}

How can this code be improved?

\begin{Shaded}
\begin{Highlighting}[]
\ControlFlowTok{if}\NormalTok{ balance }\OperatorTok{\textless{}} \DecValTok{0}\NormalTok{ :}
\NormalTok{    status }\OperatorTok{=} \StringTok{"Overdrawn"}
    \BuiltInTok{print}\NormalTok{(status)}
\ControlFlowTok{else}\NormalTok{ :}
\NormalTok{    status }\OperatorTok{=} \StringTok{"OK"}
    \BuiltInTok{print}\NormalTok{(status)}
\end{Highlighting}
\end{Shaded}

\textbf{Answer}

Move the duplicated \texttt{print} statement outside the branches.

\begin{Shaded}
\begin{Highlighting}[]
\ControlFlowTok{if}\NormalTok{ balance }\OperatorTok{\textless{}} \DecValTok{0}\NormalTok{ :}
\NormalTok{    status }\OperatorTok{=} \StringTok{"Overdrawn"}
\ControlFlowTok{else}\NormalTok{ :}
\NormalTok{    status }\OperatorTok{=} \StringTok{"OK"}

\BuiltInTok{print}\NormalTok{(status)}
\end{Highlighting}
\end{Shaded}

\begin{center}\rule{0.5\linewidth}{0.5pt}\end{center}

\section{3.2 Relational Operators}\label{relational-operators}

An \texttt{if} statement needs a condition whose value is either
\texttt{True} or \texttt{False}.

Many conditions compare two values using a \textbf{relational operator}.

\subsection{3.2.1 Relational Operators}\label{relational-operators-1}

\begin{longtable}[]{@{}ll@{}}
\toprule\noalign{}
Python operator & Meaning \\
\midrule\noalign{}
\endhead
\bottomrule\noalign{}
\endlastfoot
\texttt{\textgreater{}} & Greater than \\
\texttt{\textgreater{}=} & Greater than or equal to \\
\texttt{\textless{}} & Less than \\
\texttt{\textless{}=} & Less than or equal to \\
\texttt{==} & Equal to \\
\texttt{!=} & Not equal to \\
\end{longtable}

Examples:

\begin{Shaded}
\begin{Highlighting}[]
\DecValTok{3} \OperatorTok{\textless{}} \DecValTok{4}
\DecValTok{4} \OperatorTok{\textless{}=} \DecValTok{4}
\DecValTok{7} \OperatorTok{!=} \DecValTok{5}
\DecValTok{10} \OperatorTok{==} \DecValTok{2} \OperatorTok{*} \DecValTok{5}
\end{Highlighting}
\end{Shaded}

Each expression evaluates to a Boolean value.

\begin{center}\rule{0.5\linewidth}{0.5pt}\end{center}

\subsection{\texorpdfstring{3.2.2 Assignment \texttt{=} versus Equality
\texttt{==}}{3.2.2 Assignment = versus Equality ==}}\label{assignment-versus-equality}

These operators have very different purposes.

Assignment:

\begin{Shaded}
\begin{Highlighting}[]
\NormalTok{floor }\OperatorTok{=} \DecValTok{13}
\end{Highlighting}
\end{Shaded}

Equality test:

\begin{Shaded}
\begin{Highlighting}[]
\ControlFlowTok{if}\NormalTok{ floor }\OperatorTok{==} \DecValTok{13}\NormalTok{ :}
    \BuiltInTok{print}\NormalTok{(}\StringTok{"Floor 13 selected"}\NormalTok{)}
\end{Highlighting}
\end{Shaded}

A useful reading rule is:

\begin{Shaded}
\begin{Highlighting}[]
\NormalTok{=   → assign}
\NormalTok{==  → is equal to?}
\end{Highlighting}
\end{Shaded}

\begin{center}\rule{0.5\linewidth}{0.5pt}\end{center}

\subsection{3.2.3 Comparing Strings}\label{comparing-strings}

Strings can also be compared.

\begin{Shaded}
\begin{Highlighting}[]
\NormalTok{name1 }\OperatorTok{=} \StringTok{"John Wayne"}
\NormalTok{name2 }\OperatorTok{=} \StringTok{"John Wayne"}

\ControlFlowTok{if}\NormalTok{ name1 }\OperatorTok{==}\NormalTok{ name2 :}
    \BuiltInTok{print}\NormalTok{(}\StringTok{"The strings are identical."}\NormalTok{)}
\end{Highlighting}
\end{Shaded}

String comparison is case-sensitive:

\begin{Shaded}
\begin{Highlighting}[]
\CommentTok{"John"} \OperatorTok{==} \StringTok{"john"}
\end{Highlighting}
\end{Shaded}

is \texttt{False}.

For equality, every character and its position must match.

\begin{center}\rule{0.5\linewidth}{0.5pt}\end{center}

\subsection{3.2.4 Operator Precedence}\label{operator-precedence}

Arithmetic operators have higher precedence than relational operators.

\begin{Shaded}
\begin{Highlighting}[]
\NormalTok{floor }\OperatorTok{{-}} \DecValTok{1} \OperatorTok{\textless{}} \DecValTok{13}
\end{Highlighting}
\end{Shaded}

Python first computes:

\begin{Shaded}
\begin{Highlighting}[]
\NormalTok{floor {-} 1}
\end{Highlighting}
\end{Shaded}

and then compares the result with \texttt{13}.

For complicated expressions, parentheses can still improve readability.

\begin{center}\rule{0.5\linewidth}{0.5pt}\end{center}

\section{Common Error 3.2 -- Exact Comparison of Floating-Point
Numbers}\label{common-error-3.2-exact-comparison-of-floating-point-numbers}

Floating-point calculations may contain tiny roundoff differences.

For example:

\begin{Shaded}
\begin{Highlighting}[]
\ImportTok{from}\NormalTok{ math }\ImportTok{import}\NormalTok{ sqrt}

\NormalTok{x }\OperatorTok{=}\NormalTok{ sqrt(}\FloatTok{2.0}\NormalTok{)}
\BuiltInTok{print}\NormalTok{(x }\OperatorTok{*}\NormalTok{ x)}
\end{Highlighting}
\end{Shaded}

The result may be extremely close to \texttt{2.0} without being
represented internally as exactly \texttt{2.0}.

Therefore, this test may fail:

\begin{Shaded}
\begin{Highlighting}[]
\ControlFlowTok{if}\NormalTok{ x }\OperatorTok{*}\NormalTok{ x }\OperatorTok{==} \FloatTok{2.0}\NormalTok{ :}
    \BuiltInTok{print}\NormalTok{(}\StringTok{"Exactly two"}\NormalTok{)}
\end{Highlighting}
\end{Shaded}

When approximate equality is intended, compare the difference with a
small tolerance:

\begin{Shaded}
\begin{Highlighting}[]
\NormalTok{EPSILON }\OperatorTok{=} \FloatTok{1E{-}14}

\ControlFlowTok{if} \BuiltInTok{abs}\NormalTok{(x }\OperatorTok{*}\NormalTok{ x }\OperatorTok{{-}} \FloatTok{2.0}\NormalTok{) }\OperatorTok{\textless{}}\NormalTok{ EPSILON :}
    \BuiltInTok{print}\NormalTok{(}\StringTok{"Approximately two"}\NormalTok{)}
\end{Highlighting}
\end{Shaded}

General pattern:

\begin{Shaded}
\begin{Highlighting}[]
\ControlFlowTok{if} \BuiltInTok{abs}\NormalTok{(x }\OperatorTok{{-}}\NormalTok{ y) }\OperatorTok{\textless{}}\NormalTok{ epsilon :}
    \CommentTok{\# Treat x and y as approximately equal.}
\end{Highlighting}
\end{Shaded}

The tolerance should be chosen according to the scale and purpose of the
computation.

\begin{center}\rule{0.5\linewidth}{0.5pt}\end{center}

\section{Special Topic 3.2 -- Lexicographic Ordering of
Strings}\label{special-topic-3.2-lexicographic-ordering-of-strings}

Relational operators such as \texttt{\textless{}} and
\texttt{\textgreater{}} can be used with strings.

Python compares strings lexicographically according to character codes.

For example, capitalization matters:

\begin{Shaded}
\begin{Highlighting}[]
\CommentTok{"John"} \OperatorTok{\textless{}} \StringTok{"john"}
\end{Highlighting}
\end{Shaded}

The result depends on the character ordering, not on dictionary rules
that ignore case.

For case-insensitive comparisons, a useful technique is to normalize
both strings first:

\begin{Shaded}
\begin{Highlighting}[]
\NormalTok{name1 }\OperatorTok{=}\NormalTok{ name1.lower()}
\NormalTok{name2 }\OperatorTok{=}\NormalTok{ name2.lower()}
\end{Highlighting}
\end{Shaded}

and then compare them.

\begin{center}\rule{0.5\linewidth}{0.5pt}\end{center}

\section{\texorpdfstring{HOW TO 3.1 -- Implementing an \texttt{if}
Statement}{HOW TO 3.1 -- Implementing an if Statement}}\label{how-to-3.1-implementing-an-if-statement}

The chapter proposes a systematic approach to decisions.

Suppose a store applies one discount rate to prices below a threshold
and a larger discount at or above the threshold.

\subsection{Step 1. Decide on the branching
condition}\label{step-1.-decide-on-the-branching-condition}

Ask a question that can be answered \texttt{True} or \texttt{False}.

Example:

\begin{Shaded}
\begin{Highlighting}[]
\NormalTok{Is the original price less than 128?}
\end{Highlighting}
\end{Shaded}

\begin{center}\rule{0.5\linewidth}{0.5pt}\end{center}

\subsection{Step 2. Describe what happens when the condition is
true}\label{step-2.-describe-what-happens-when-the-condition-is-true}

\begin{Shaded}
\begin{Highlighting}[]
\NormalTok{Use the lower discount rate.}
\end{Highlighting}
\end{Shaded}

\begin{center}\rule{0.5\linewidth}{0.5pt}\end{center}

\subsection{Step 3. Describe what happens when the condition is
false}\label{step-3.-describe-what-happens-when-the-condition-is-false}

\begin{Shaded}
\begin{Highlighting}[]
\NormalTok{Use the higher discount rate.}
\end{Highlighting}
\end{Shaded}

\begin{center}\rule{0.5\linewidth}{0.5pt}\end{center}

\subsection{Step 4. Check the relational operator
carefully}\label{step-4.-check-the-relational-operator-carefully}

Ask what should happen at the boundary.

If the higher discount begins at exactly \texttt{128}, then the
lower-rate condition should be:

\begin{Shaded}
\begin{Highlighting}[]
\NormalTok{originalPrice }\OperatorTok{\textless{}} \DecValTok{128}
\end{Highlighting}
\end{Shaded}

not:

\begin{Shaded}
\begin{Highlighting}[]
\NormalTok{originalPrice }\OperatorTok{\textless{}=} \DecValTok{128}
\end{Highlighting}
\end{Shaded}

Boundary values often reveal incorrect \texttt{\textless{}} versus
\texttt{\textless{}=} choices.

\begin{center}\rule{0.5\linewidth}{0.5pt}\end{center}

\subsection{Step 5. Remove duplicated
work}\label{step-5.-remove-duplicated-work}

Instead of computing the entire discounted price in each branch, choose
only the rate inside the branches:

\begin{Shaded}
\begin{Highlighting}[]
\ControlFlowTok{if}\NormalTok{ originalPrice }\OperatorTok{\textless{}} \DecValTok{128}\NormalTok{ :}
\NormalTok{    discountRate }\OperatorTok{=} \FloatTok{0.92}
\ControlFlowTok{else}\NormalTok{ :}
\NormalTok{    discountRate }\OperatorTok{=} \FloatTok{0.84}

\NormalTok{discountedPrice }\OperatorTok{=}\NormalTok{ discountRate }\OperatorTok{*}\NormalTok{ originalPrice}
\end{Highlighting}
\end{Shaded}

\begin{center}\rule{0.5\linewidth}{0.5pt}\end{center}

\subsection{Step 6. Test both
branches}\label{step-6.-test-both-branches}

Choose at least one value below the threshold and one above it.

For example:

\begin{Shaded}
\begin{Highlighting}[]
\NormalTok{100  → lower{-}discount branch}
\NormalTok{200  → higher{-}discount branch}
\NormalTok{128  → boundary case}
\end{Highlighting}
\end{Shaded}

\begin{center}\rule{0.5\linewidth}{0.5pt}\end{center}

\subsection{Step 7. Assemble the Python
program}\label{step-7.-assemble-the-python-program}

\begin{Shaded}
\begin{Highlighting}[]
\NormalTok{originalPrice }\OperatorTok{=} \BuiltInTok{float}\NormalTok{(}\BuiltInTok{input}\NormalTok{(}\StringTok{"Original price before discount: "}\NormalTok{))}

\ControlFlowTok{if}\NormalTok{ originalPrice }\OperatorTok{\textless{}} \DecValTok{128}\NormalTok{ :}
\NormalTok{    discountRate }\OperatorTok{=} \FloatTok{0.92}
\ControlFlowTok{else}\NormalTok{ :}
\NormalTok{    discountRate }\OperatorTok{=} \FloatTok{0.84}

\NormalTok{discountedPrice }\OperatorTok{=}\NormalTok{ discountRate }\OperatorTok{*}\NormalTok{ originalPrice}
\BuiltInTok{print}\NormalTok{(}\StringTok{"Discounted price: }\SpecialCharTok{\%.2f}\StringTok{"} \OperatorTok{\%}\NormalTok{ discountedPrice)}
\end{Highlighting}
\end{Shaded}

The important skill is not memorizing this particular price example. It
is learning how to translate a two-case rule into a correct condition
and two branches.

\begin{center}\rule{0.5\linewidth}{0.5pt}\end{center}

\section{Worked Example 3.1 -- Extracting the Middle of a
String}\label{worked-example-3.1-extracting-the-middle-of-a-string}

Suppose a program must extract:

\begin{itemize}
\tightlist
\item
  one middle character from an odd-length string;
\item
  two middle characters from an even-length string.
\end{itemize}

Examples:

\begin{Shaded}
\begin{Highlighting}[]
\NormalTok{"crate"   → "a"}
\NormalTok{"crates"  → "at"}
\end{Highlighting}
\end{Shaded}

The branching condition is based on parity:

\begin{Shaded}
\begin{Highlighting}[]
\BuiltInTok{len}\NormalTok{(text) }\OperatorTok{\%} \DecValTok{2} \OperatorTok{==} \DecValTok{1}
\end{Highlighting}
\end{Shaded}

The middle position can be computed once:

\begin{Shaded}
\begin{Highlighting}[]
\NormalTok{position }\OperatorTok{=} \BuiltInTok{len}\NormalTok{(text) }\OperatorTok{//} \DecValTok{2}
\end{Highlighting}
\end{Shaded}

Then:

\begin{Shaded}
\begin{Highlighting}[]
\ControlFlowTok{if} \BuiltInTok{len}\NormalTok{(text) }\OperatorTok{\%} \DecValTok{2} \OperatorTok{==} \DecValTok{1}\NormalTok{ :}
\NormalTok{    result }\OperatorTok{=}\NormalTok{ text[position]}
\ControlFlowTok{else}\NormalTok{ :}
\NormalTok{    result }\OperatorTok{=}\NormalTok{ text[position }\OperatorTok{{-}} \DecValTok{1}\NormalTok{] }\OperatorTok{+}\NormalTok{ text[position]}
\end{Highlighting}
\end{Shaded}

This example combines:

\begin{itemize}
\tightlist
\item
  string length;
\item
  floor division;
\item
  remainder;
\item
  indexing;
\item
  a two-way decision;
\item
  removal of duplicated calculations.
\end{itemize}

\begin{center}\rule{0.5\linewidth}{0.5pt}\end{center}

\subsection{Self Check 3.2}\label{self-check-3.2}

\subsubsection{Question 1}\label{question-1-1}

What is the value of:

\begin{Shaded}
\begin{Highlighting}[]
\DecValTok{4} \OperatorTok{\textless{}=} \DecValTok{4}
\end{Highlighting}
\end{Shaded}

\textbf{Answer}

\texttt{True}.

\subsubsection{Question 2}\label{question-2-1}

Why is this incorrect?

\begin{Shaded}
\begin{Highlighting}[]
\ControlFlowTok{if}\NormalTok{ score }\OperatorTok{=} \DecValTok{100}\NormalTok{ :}
    \BuiltInTok{print}\NormalTok{(}\StringTok{"Perfect"}\NormalTok{)}
\end{Highlighting}
\end{Shaded}

\textbf{Answer}

\texttt{=} performs assignment. Equality testing requires \texttt{==}.

\subsubsection{Question 3}\label{question-3-1}

What condition tests whether \texttt{n} is odd?

\textbf{Answer}

\begin{Shaded}
\begin{Highlighting}[]
\NormalTok{n }\OperatorTok{\%} \DecValTok{2} \OperatorTok{==} \DecValTok{1}
\end{Highlighting}
\end{Shaded}

\subsubsection{Question 4}\label{question-4-1}

What is the middle result for:

\begin{Shaded}
\begin{Highlighting}[]
\NormalTok{text }\OperatorTok{=} \StringTok{"monitor"}
\end{Highlighting}
\end{Shaded}

\textbf{Answer}

\texttt{"i"}.

\begin{center}\rule{0.5\linewidth}{0.5pt}\end{center}

\section{3.3 Nested Branches}\label{nested-branches}

Sometimes a program must make one decision and then make another
decision inside the selected branch.

This is called \textbf{nesting}.

Example structure:

\begin{Shaded}
\begin{Highlighting}[]
\ControlFlowTok{if}\NormalTok{ condition1 :}
    \ControlFlowTok{if}\NormalTok{ condition2 :}
\NormalTok{        statementA}
    \ControlFlowTok{else}\NormalTok{ :}
\NormalTok{        statementB}
\ControlFlowTok{else}\NormalTok{ :}
\NormalTok{    statementC}
\end{Highlighting}
\end{Shaded}

The second \texttt{if} is part of the first branch.

\begin{center}\rule{0.5\linewidth}{0.5pt}\end{center}

\subsection{3.3.1 Multi-Level Decisions}\label{multi-level-decisions}

A useful example is a simplified tax calculation in which:

\begin{enumerate}
\def\labelenumi{\arabic{enumi}.}
\tightlist
\item
  the program first determines a taxpayer category;
\item
  it then determines which income bracket applies within that category.
\end{enumerate}

Conceptually:

\begin{Shaded}
\begin{Highlighting}[]
\NormalTok{                     status?}
\NormalTok{                 /             \textbackslash{}}
\NormalTok{             single           married}
\NormalTok{              /                  \textbackslash{}}
\NormalTok{        income limit?        income limit?}
\NormalTok{          /      \textbackslash{}             /      \textbackslash{}}
\NormalTok{      lower     upper       lower     upper}
\NormalTok{      rate      rate        rate      rate}
\end{Highlighting}
\end{Shaded}

The important point is the structure: \textbf{one decision depends on
the result of another decision}.

A simplified implementation pattern is:

\begin{Shaded}
\begin{Highlighting}[]
\ControlFlowTok{if}\NormalTok{ status }\OperatorTok{==} \StringTok{"s"}\NormalTok{ :}
    \ControlFlowTok{if}\NormalTok{ income }\OperatorTok{\textless{}=}\NormalTok{ SINGLE\_LIMIT :}
        \CommentTok{\# lower bracket for single}
\NormalTok{        ...}
    \ControlFlowTok{else}\NormalTok{ :}
        \CommentTok{\# upper bracket for single}
\NormalTok{        ...}
\ControlFlowTok{else}\NormalTok{ :}
    \ControlFlowTok{if}\NormalTok{ income }\OperatorTok{\textless{}=}\NormalTok{ MARRIED\_LIMIT :}
        \CommentTok{\# lower bracket for married}
\NormalTok{        ...}
    \ControlFlowTok{else}\NormalTok{ :}
        \CommentTok{\# upper bracket for married}
\NormalTok{        ...}
\end{Highlighting}
\end{Shaded}

Nested decisions can go deeper, but excessive nesting makes programs
harder to read. Later lessons will introduce functions as one way to
organize more complex logic.

\begin{center}\rule{0.5\linewidth}{0.5pt}\end{center}

\section{Programming Tip 3.2 --
Hand-Tracing}\label{programming-tip-3.2-hand-tracing}

\textbf{Hand-tracing} means simulating the program manually, one
statement at a time.

A useful method is to create a small table:

\begin{longtable}[]{@{}lrlrr@{}}
\toprule\noalign{}
Step & \texttt{income} & \texttt{status} & \texttt{tax1} &
\texttt{tax2} \\
\midrule\noalign{}
\endhead
\bottomrule\noalign{}
\endlastfoot
Input & 80000 & \texttt{"m"} & & \\
Initialize & 80000 & \texttt{"m"} & 0 & 0 \\
Branch & 80000 & \texttt{"m"} & \ldots{} & \ldots{} \\
\end{longtable}

As each assignment occurs, update the affected variable.

Hand-tracing is especially useful for:

\begin{itemize}
\tightlist
\item
  nested \texttt{if} statements;
\item
  code with several variables;
\item
  understanding why a program entered a particular branch;
\item
  comparing expected and actual program behavior.
\end{itemize}

\begin{center}\rule{0.5\linewidth}{0.5pt}\end{center}

\section{Computing \& Society 3.1 -- Complex Systems and Decision
Logic}\label{computing-society-3.1-complex-systems-and-decision-logic}

The chapter uses the Denver International Airport baggage-handling
system as a case study in software and system complexity. The broader
lesson is that decision logic that appears manageable in isolation can
become difficult when many physical components, timing constraints,
exceptional cases, and interactions are combined.

For beginning programmers, the practical lesson is:

\begin{quote}
Keep each decision understandable, test the branches systematically, and
do not assume that a large system is simply a collection of
independently correct small pieces.
\end{quote}

\begin{center}\rule{0.5\linewidth}{0.5pt}\end{center}

\subsection{Self Check 3.3}\label{self-check-3.3}

\subsubsection{Question 1}\label{question-1-2}

When is nesting appropriate?

\textbf{Answer}

When a second decision only makes sense after a particular result of an
earlier decision.

\subsubsection{Question 2}\label{question-2-2}

Why is indentation particularly important in nested decisions?

\textbf{Answer}

Indentation determines which outer branch contains each inner statement.

\subsubsection{Question 3}\label{question-3-2}

What does hand-tracing help you understand?

\textbf{Answer}

The order of execution, which branches are selected, and how variable
values change.

\begin{center}\rule{0.5\linewidth}{0.5pt}\end{center}

\section{3.4 Multiple Alternatives}\label{multiple-alternatives}

Many decisions have more than two possible outcomes.

Python uses \texttt{elif} to express a sequence of mutually exclusive
alternatives.

General form:

\begin{Shaded}
\begin{Highlighting}[]
\ControlFlowTok{if}\NormalTok{ condition1 :}
\NormalTok{    statements1}
\ControlFlowTok{elif}\NormalTok{ condition2 :}
\NormalTok{    statements2}
\ControlFlowTok{elif}\NormalTok{ condition3 :}
\NormalTok{    statements3}
\ControlFlowTok{else}\NormalTok{ :}
\NormalTok{    defaultStatements}
\end{Highlighting}
\end{Shaded}

Python evaluates the conditions from top to bottom. As soon as one
condition is true, its branch executes and the remaining alternatives
are skipped.

\begin{center}\rule{0.5\linewidth}{0.5pt}\end{center}

\subsection{Example -- Classifying an
Earthquake}\label{example-classifying-an-earthquake}

A simplified classification may use thresholds such as:

\begin{Shaded}
\begin{Highlighting}[]
\NormalTok{richter }\OperatorTok{=} \BuiltInTok{float}\NormalTok{(}\BuiltInTok{input}\NormalTok{(}\StringTok{"Enter a magnitude: "}\NormalTok{))}

\ControlFlowTok{if}\NormalTok{ richter }\OperatorTok{\textgreater{}=} \FloatTok{8.0}\NormalTok{ :}
\NormalTok{    description }\OperatorTok{=} \StringTok{"Very severe structural damage"}
\ControlFlowTok{elif}\NormalTok{ richter }\OperatorTok{\textgreater{}=} \FloatTok{7.0}\NormalTok{ :}
\NormalTok{    description }\OperatorTok{=} \StringTok{"Severe damage in many areas"}
\ControlFlowTok{elif}\NormalTok{ richter }\OperatorTok{\textgreater{}=} \FloatTok{6.0}\NormalTok{ :}
\NormalTok{    description }\OperatorTok{=} \StringTok{"Considerable damage is possible"}
\ControlFlowTok{elif}\NormalTok{ richter }\OperatorTok{\textgreater{}=} \FloatTok{4.5}\NormalTok{ :}
\NormalTok{    description }\OperatorTok{=} \StringTok{"Some poorly constructed buildings may be damaged"}
\ControlFlowTok{else}\NormalTok{ :}
\NormalTok{    description }\OperatorTok{=} \StringTok{"Little or no structural damage"}

\BuiltInTok{print}\NormalTok{(description)}
\end{Highlighting}
\end{Shaded}

The exact descriptive labels are less important than the branching
structure.

\begin{center}\rule{0.5\linewidth}{0.5pt}\end{center}

\subsection{3.4.1 Order Matters}\label{order-matters}

Suppose the conditions are written from the smallest threshold upward:

\begin{Shaded}
\begin{Highlighting}[]
\ControlFlowTok{if}\NormalTok{ richter }\OperatorTok{\textgreater{}=} \FloatTok{4.5}\NormalTok{ :}
\NormalTok{    ...}
\ControlFlowTok{elif}\NormalTok{ richter }\OperatorTok{\textgreater{}=} \FloatTok{6.0}\NormalTok{ :}
\NormalTok{    ...}
\ControlFlowTok{elif}\NormalTok{ richter }\OperatorTok{\textgreater{}=} \FloatTok{7.0}\NormalTok{ :}
\NormalTok{    ...}
\end{Highlighting}
\end{Shaded}

For \texttt{richter\ =\ 7.1}, the first condition is already true, so
later branches are never considered.

When ranges overlap, test the \textbf{more specific / higher threshold}
first.

Correct pattern:

\begin{Shaded}
\begin{Highlighting}[]
\ControlFlowTok{if}\NormalTok{ value }\OperatorTok{\textgreater{}=} \DecValTok{80}\NormalTok{ :}
\NormalTok{    ...}
\ControlFlowTok{elif}\NormalTok{ value }\OperatorTok{\textgreater{}=} \DecValTok{70}\NormalTok{ :}
\NormalTok{    ...}
\ControlFlowTok{elif}\NormalTok{ value }\OperatorTok{\textgreater{}=} \DecValTok{60}\NormalTok{ :}
\NormalTok{    ...}
\ControlFlowTok{else}\NormalTok{ :}
\NormalTok{    ...}
\end{Highlighting}
\end{Shaded}

\begin{center}\rule{0.5\linewidth}{0.5pt}\end{center}

\subsection{\texorpdfstring{3.4.2 \texttt{elif} versus Independent
\texttt{if}
Statements}{3.4.2 elif versus Independent if Statements}}\label{elif-versus-independent-if-statements}

These are not equivalent.

Mutually exclusive alternatives:

\begin{Shaded}
\begin{Highlighting}[]
\ControlFlowTok{if}\NormalTok{ score }\OperatorTok{\textgreater{}=} \DecValTok{90}\NormalTok{ :}
\NormalTok{    grade }\OperatorTok{=} \StringTok{"A"}
\ControlFlowTok{elif}\NormalTok{ score }\OperatorTok{\textgreater{}=} \DecValTok{80}\NormalTok{ :}
\NormalTok{    grade }\OperatorTok{=} \StringTok{"B"}
\ControlFlowTok{elif}\NormalTok{ score }\OperatorTok{\textgreater{}=} \DecValTok{70}\NormalTok{ :}
\NormalTok{    grade }\OperatorTok{=} \StringTok{"C"}
\end{Highlighting}
\end{Shaded}

Only one branch executes.

Independent tests:

\begin{Shaded}
\begin{Highlighting}[]
\ControlFlowTok{if}\NormalTok{ score }\OperatorTok{\textgreater{}=} \DecValTok{90}\NormalTok{ :}
    \BuiltInTok{print}\NormalTok{(}\StringTok{"At least 90"}\NormalTok{)}
\ControlFlowTok{if}\NormalTok{ score }\OperatorTok{\textgreater{}=} \DecValTok{80}\NormalTok{ :}
    \BuiltInTok{print}\NormalTok{(}\StringTok{"At least 80"}\NormalTok{)}
\ControlFlowTok{if}\NormalTok{ score }\OperatorTok{\textgreater{}=} \DecValTok{70}\NormalTok{ :}
    \BuiltInTok{print}\NormalTok{(}\StringTok{"At least 70"}\NormalTok{)}
\end{Highlighting}
\end{Shaded}

For \texttt{score\ =\ 95}, all three statements print.

Use \texttt{if/elif/else} when alternatives are intended to be mutually
exclusive. Use separate \texttt{if} statements when several conditions
may independently trigger actions.

\begin{center}\rule{0.5\linewidth}{0.5pt}\end{center}

\subsection{Self Check 3.4}\label{self-check-3.4}

\subsubsection{Question 1}\label{question-1-3}

Write the logic for setting \texttt{sign} to \texttt{1}, \texttt{0}, or
\texttt{-1} depending on whether \texttt{x} is positive, zero, or
negative.

\textbf{Answer}

\begin{Shaded}
\begin{Highlighting}[]
\ControlFlowTok{if}\NormalTok{ x }\OperatorTok{\textgreater{}} \DecValTok{0}\NormalTok{ :}
\NormalTok{    sign }\OperatorTok{=} \DecValTok{1}
\ControlFlowTok{elif}\NormalTok{ x }\OperatorTok{\textless{}} \DecValTok{0}\NormalTok{ :}
\NormalTok{    sign }\OperatorTok{=} \OperatorTok{{-}}\DecValTok{1}
\ControlFlowTok{else}\NormalTok{ :}
\NormalTok{    sign }\OperatorTok{=} \DecValTok{0}
\end{Highlighting}
\end{Shaded}

\subsubsection{Question 2}\label{question-2-3}

Why should overlapping threshold conditions usually be ordered from
largest to smallest?

\textbf{Answer}

Because the first true branch ends the \texttt{if/elif} chain. A broad
low-threshold condition placed first can capture values that should
belong to a later, more specific case.

\begin{center}\rule{0.5\linewidth}{0.5pt}\end{center}

\section{TOOLBOX 3.1 -- Sending
E-mail}\label{toolbox-3.1-sending-e-mail}

The textbook uses e-mail automation as an example of combining data,
decisions, and a Python library.

The key programming idea is to construct different message content
according to a condition:

\begin{Shaded}
\begin{Highlighting}[]
\NormalTok{score }\OperatorTok{=} \BuiltInTok{int}\NormalTok{(}\BuiltInTok{input}\NormalTok{(}\StringTok{"Score: "}\NormalTok{))}
\NormalTok{body }\OperatorTok{=} \StringTok{"Your score on the last exam is "} \OperatorTok{+} \BuiltInTok{str}\NormalTok{(score) }\OperatorTok{+} \StringTok{"}\CharTok{\textbackslash{}n}\StringTok{"}

\ControlFlowTok{if}\NormalTok{ score }\OperatorTok{\textless{}=} \DecValTok{50}\NormalTok{ :}
\NormalTok{    body }\OperatorTok{+=} \StringTok{"Please consider visiting the tutoring center."}
\ControlFlowTok{elif}\NormalTok{ score }\OperatorTok{\textgreater{}=} \DecValTok{90}\NormalTok{ :}
\NormalTok{    body }\OperatorTok{+=} \StringTok{"Excellent work."}
\end{Highlighting}
\end{Shaded}

The library-specific details are optional for this lesson. The
decision-making concept is that the program can customize an action or
message according to data.

\begin{quote}
For real accounts, authentication and mail-provider requirements change
over time. Do not place real passwords directly in teaching code or
shared notebooks.
\end{quote}

\begin{center}\rule{0.5\linewidth}{0.5pt}\end{center}

\section{3.5 Problem Solving:
Flowcharts}\label{problem-solving-flowcharts}

A \textbf{flowchart} visualizes the flow of control in an algorithm.

Common elements include:

\begin{Shaded}
\begin{Highlighting}[]
\NormalTok{[ Process / task ]}

\NormalTok{\textless{} Input / output \textgreater{}}

\NormalTok{      / Condition? \textbackslash{}}
\NormalTok{     /             \textbackslash{}}
\NormalTok{  True             False}
\end{Highlighting}
\end{Shaded}

Flowcharts help when a problem contains several decisions and it is not
yet clear how the branches should fit together.

\begin{center}\rule{0.5\linewidth}{0.5pt}\end{center}

\subsection{3.5.1 Two-Way Decision}\label{two-way-decision}

\begin{Shaded}
\begin{Highlighting}[]
\NormalTok{              temperature \textless{} 0 ?}
\NormalTok{                  /        \textbackslash{}}
\NormalTok{              True        False}
\NormalTok{               |             |}
\NormalTok{        print("Frozen")    continue}
\end{Highlighting}
\end{Shaded}

\begin{center}\rule{0.5\linewidth}{0.5pt}\end{center}

\subsection{3.5.2 Multiple Alternatives}\label{multiple-alternatives-1}

\begin{Shaded}
\begin{Highlighting}[]
\NormalTok{                score \textgreater{}= 90 ?}
\NormalTok{                 /       \textbackslash{}}
\NormalTok{              True       False}
\NormalTok{               |           |}
\NormalTok{           grade A      score \textgreater{}= 80 ?}
\NormalTok{                         /       \textbackslash{}}
\NormalTok{                      True       False}
\NormalTok{                       |           |}
\NormalTok{                    grade B      ...}
\end{Highlighting}
\end{Shaded}

\begin{center}\rule{0.5\linewidth}{0.5pt}\end{center}

\subsection{3.5.3 Avoiding Spaghetti
Code}\label{avoiding-spaghetti-code}

A useful design rule from the chapter is:

\begin{quote}
Do not draw an arrow from one branch into the middle of another branch.
\end{quote}

Branches should be structured so that they clearly split and later
rejoin.

This keeps the control flow consistent with structured \texttt{if},
\texttt{if/else}, and nested statements.

\begin{center}\rule{0.5\linewidth}{0.5pt}\end{center}

\subsection{Example -- Shipping Cost}\label{example-shipping-cost}

Suppose:

\begin{itemize}
\tightlist
\item
  shipping inside the continental USA costs one rate;
\item
  shipping to Alaska or Hawaii costs another rate;
\item
  international shipping uses the higher rate.
\end{itemize}

A structured solution is:

\begin{Shaded}
\begin{Highlighting}[]
\NormalTok{country }\OperatorTok{=} \BuiltInTok{input}\NormalTok{(}\StringTok{"Enter the country: "}\NormalTok{)}
\NormalTok{state }\OperatorTok{=} \BuiltInTok{input}\NormalTok{(}\StringTok{"Enter the state or province: "}\NormalTok{)}

\ControlFlowTok{if}\NormalTok{ country }\OperatorTok{==} \StringTok{"USA"}\NormalTok{ :}
    \ControlFlowTok{if}\NormalTok{ state }\OperatorTok{==} \StringTok{"AK"} \KeywordTok{or}\NormalTok{ state }\OperatorTok{==} \StringTok{"HI"}\NormalTok{ :}
\NormalTok{        shippingCost }\OperatorTok{=} \FloatTok{10.0}
    \ControlFlowTok{else}\NormalTok{ :}
\NormalTok{        shippingCost }\OperatorTok{=} \FloatTok{5.0}
\ControlFlowTok{else}\NormalTok{ :}
\NormalTok{    shippingCost }\OperatorTok{=} \FloatTok{10.0}

\BuiltInTok{print}\NormalTok{(}\StringTok{"Shipping cost to }\SpecialCharTok{\%s}\StringTok{, }\SpecialCharTok{\%s}\StringTok{: $}\SpecialCharTok{\%.2f}\StringTok{"} \OperatorTok{\%}
\NormalTok{      (state, country, shippingCost))}
\end{Highlighting}
\end{Shaded}

The same structure can first be drawn as a flowchart and then translated
into Python.

\begin{center}\rule{0.5\linewidth}{0.5pt}\end{center}

\subsection{Self Check 3.5}\label{self-check-3.5}

\subsubsection{Question 1}\label{question-1-4}

What is the main purpose of a flowchart?

\textbf{Answer}

To visualize tasks, decisions, and the possible paths of execution
before or while implementing the algorithm.

\subsubsection{Question 2}\label{question-2-4}

Why should arrows not enter the middle of another branch?

\textbf{Answer}

It creates unstructured control flow that is difficult to translate,
test, and maintain.

\begin{center}\rule{0.5\linewidth}{0.5pt}\end{center}

\section{3.6 Problem Solving: Test
Cases}\label{problem-solving-test-cases}

A program with decisions must be tested with inputs that exercise the
different paths through the program.

Trying one ordinary input is not enough.

A strong test plan aims for:

\begin{enumerate}
\def\labelenumi{\arabic{enumi}.}
\tightlist
\item
  \textbf{branch coverage} -- each branch is executed by at least one
  test;
\item
  \textbf{boundary testing} -- values exactly at or very near decision
  boundaries are tested;
\item
  \textbf{invalid-input testing} -- when the program is responsible for
  validating input.
\end{enumerate}

\begin{center}\rule{0.5\linewidth}{0.5pt}\end{center}

\subsection{3.6.1 Branch Coverage}\label{branch-coverage}

For a two-way decision:

\begin{Shaded}
\begin{Highlighting}[]
\ControlFlowTok{if}\NormalTok{ age }\OperatorTok{\textgreater{}=} \DecValTok{18}\NormalTok{ :}
\NormalTok{    category }\OperatorTok{=} \StringTok{"Adult"}
\ControlFlowTok{else}\NormalTok{ :}
\NormalTok{    category }\OperatorTok{=} \StringTok{"Minor"}
\end{Highlighting}
\end{Shaded}

Use at least one test for each outcome:

\begin{longtable}[]{@{}rll@{}}
\toprule\noalign{}
Input & Expected result & Purpose \\
\midrule\noalign{}
\endhead
\bottomrule\noalign{}
\endlastfoot
\texttt{20} & Adult & True branch \\
\texttt{15} & Minor & False branch \\
\end{longtable}

\begin{center}\rule{0.5\linewidth}{0.5pt}\end{center}

\subsection{3.6.2 Boundary Cases}\label{boundary-cases}

The boundary is \texttt{18}.

Useful tests include:

\begin{Shaded}
\begin{Highlighting}[]
\NormalTok{17}
\NormalTok{18}
\NormalTok{19}
\end{Highlighting}
\end{Shaded}

The value \texttt{18} is particularly important because it distinguishes
\texttt{\textgreater{}=} from \texttt{\textgreater{}}.

\begin{center}\rule{0.5\linewidth}{0.5pt}\end{center}

\subsection{3.6.3 Design Tests Before
Coding}\label{design-tests-before-coding}

Writing expected results before implementation has several benefits:

\begin{itemize}
\tightlist
\item
  forces you to understand the rule;
\item
  reveals ambiguous boundary conditions;
\item
  gives you an immediate way to verify the finished program;
\item
  reduces the temptation to accept whatever output the program happens
  to produce.
\end{itemize}

A useful test table is:

\begin{longtable}[]{@{}lll@{}}
\toprule\noalign{}
Test case & Expected output & Reason \\
\midrule\noalign{}
\endhead
\bottomrule\noalign{}
\endlastfoot
ordinary case A & \ldots{} & first branch \\
ordinary case B & \ldots{} & second branch \\
boundary & \ldots{} & exact cutoff \\
invalid value & error & validation \\
\end{longtable}

\begin{center}\rule{0.5\linewidth}{0.5pt}\end{center}

\section{Programming Tip 3.3 -- Plan Time for Testing and
Debugging}\label{programming-tip-3.3-plan-time-for-testing-and-debugging}

Programming work includes more than typing code. A realistic plan
includes time for:

\begin{itemize}
\tightlist
\item
  understanding and designing the solution;
\item
  preparing test cases;
\item
  entering the program and correcting syntax errors;
\item
  running tests and debugging unexpected behavior.
\end{itemize}

Programs with more branches have more possible execution paths and
therefore usually require more testing effort.

\begin{center}\rule{0.5\linewidth}{0.5pt}\end{center}

\subsection{Self Check 3.6}\label{self-check-3.6}

\subsubsection{Question 1}\label{question-1-5}

For the condition:

\begin{Shaded}
\begin{Highlighting}[]
\ControlFlowTok{if}\NormalTok{ value }\OperatorTok{\textless{}} \DecValTok{100}\NormalTok{ :}
\end{Highlighting}
\end{Shaded}

what is an important boundary test?

\textbf{Answer}

\texttt{100}. Values such as \texttt{99} and \texttt{101} are also
useful neighbors.

\subsubsection{Question 2}\label{question-2-5}

Why is one test case usually insufficient for an \texttt{if/else}
statement?

\textbf{Answer}

It may execute only one branch, leaving the other branch untested.

\begin{center}\rule{0.5\linewidth}{0.5pt}\end{center}

\section{3.7 Boolean Variables and
Operators}\label{boolean-variables-and-operators}

The Boolean type \texttt{bool} has exactly two values:

\begin{Shaded}
\begin{Highlighting}[]
\VariableTok{True}
\VariableTok{False}
\end{Highlighting}
\end{Shaded}

A Boolean variable can store the result of a logical condition.

\begin{Shaded}
\begin{Highlighting}[]
\NormalTok{failed }\OperatorTok{=} \VariableTok{True}

\ControlFlowTok{if}\NormalTok{ failed :}
    \BuiltInTok{print}\NormalTok{(}\StringTok{"The operation failed"}\NormalTok{)}
\end{Highlighting}
\end{Shaded}

A Boolean variable is sometimes called a \textbf{flag} because its state
indicates whether a condition holds.

\begin{center}\rule{0.5\linewidth}{0.5pt}\end{center}

\subsection{\texorpdfstring{3.7.1 The \texttt{and}
Operator}{3.7.1 The and Operator}}\label{the-and-operator}

\texttt{and} is true only when \textbf{both} operands are true.

\begin{Shaded}
\begin{Highlighting}[]
\ControlFlowTok{if}\NormalTok{ temp }\OperatorTok{\textgreater{}} \DecValTok{0} \KeywordTok{and}\NormalTok{ temp }\OperatorTok{\textless{}} \DecValTok{100}\NormalTok{ :}
    \BuiltInTok{print}\NormalTok{(}\StringTok{"Liquid water"}\NormalTok{)}
\end{Highlighting}
\end{Shaded}

Truth table:

\begin{longtable}[]{@{}lll@{}}
\toprule\noalign{}
A & B & \texttt{A\ and\ B} \\
\midrule\noalign{}
\endhead
\bottomrule\noalign{}
\endlastfoot
\texttt{True} & \texttt{True} & \texttt{True} \\
\texttt{True} & \texttt{False} & \texttt{False} \\
\texttt{False} & \texttt{True} & \texttt{False} \\
\texttt{False} & \texttt{False} & \texttt{False} \\
\end{longtable}

\begin{center}\rule{0.5\linewidth}{0.5pt}\end{center}

\subsection{\texorpdfstring{3.7.2 The \texttt{or}
Operator}{3.7.2 The or Operator}}\label{the-or-operator}

\texttt{or} is true when \textbf{at least one} operand is true.

\begin{Shaded}
\begin{Highlighting}[]
\ControlFlowTok{if}\NormalTok{ temp }\OperatorTok{\textless{}=} \DecValTok{0} \KeywordTok{or}\NormalTok{ temp }\OperatorTok{\textgreater{}=} \DecValTok{100}\NormalTok{ :}
    \BuiltInTok{print}\NormalTok{(}\StringTok{"Not liquid water"}\NormalTok{)}
\end{Highlighting}
\end{Shaded}

Truth table:

\begin{longtable}[]{@{}lll@{}}
\toprule\noalign{}
A & B & \texttt{A\ or\ B} \\
\midrule\noalign{}
\endhead
\bottomrule\noalign{}
\endlastfoot
\texttt{True} & \texttt{True} & \texttt{True} \\
\texttt{True} & \texttt{False} & \texttt{True} \\
\texttt{False} & \texttt{True} & \texttt{True} \\
\texttt{False} & \texttt{False} & \texttt{False} \\
\end{longtable}

Notice that \texttt{or} is inclusive: it is still true when
\textbf{both} conditions are true.

\begin{center}\rule{0.5\linewidth}{0.5pt}\end{center}

\subsection{\texorpdfstring{3.7.3 The \texttt{not}
Operator}{3.7.3 The not Operator}}\label{the-not-operator}

\texttt{not} reverses a Boolean value.

\begin{Shaded}
\begin{Highlighting}[]
\ControlFlowTok{if} \KeywordTok{not}\NormalTok{ frozen :}
    \BuiltInTok{print}\NormalTok{(}\StringTok{"Not frozen"}\NormalTok{)}
\end{Highlighting}
\end{Shaded}

\begin{longtable}[]{@{}ll@{}}
\toprule\noalign{}
A & \texttt{not\ A} \\
\midrule\noalign{}
\endhead
\bottomrule\noalign{}
\endlastfoot
\texttt{True} & \texttt{False} \\
\texttt{False} & \texttt{True} \\
\end{longtable}

\begin{center}\rule{0.5\linewidth}{0.5pt}\end{center}

\subsection{3.7.4 Precedence}\label{precedence}

A useful simplified order is:

\begin{Shaded}
\begin{Highlighting}[]
\NormalTok{arithmetic}
\NormalTok{   ↓}
\NormalTok{relational operators}
\NormalTok{   ↓}
\NormalTok{not}
\NormalTok{   ↓}
\NormalTok{and}
\NormalTok{   ↓}
\NormalTok{or}
\end{Highlighting}
\end{Shaded}

Thus:

\begin{Shaded}
\begin{Highlighting}[]
\NormalTok{x }\OperatorTok{\textgreater{}} \DecValTok{0} \KeywordTok{and}\NormalTok{ y }\OperatorTok{\textless{}} \DecValTok{10}
\end{Highlighting}
\end{Shaded}

is interpreted as:

\begin{Shaded}
\begin{Highlighting}[]
\NormalTok{(x }\OperatorTok{\textgreater{}} \DecValTok{0}\NormalTok{) }\KeywordTok{and}\NormalTok{ (y }\OperatorTok{\textless{}} \DecValTok{10}\NormalTok{)}
\end{Highlighting}
\end{Shaded}

For complex expressions, parentheses improve readability even when they
are not required by precedence.

\begin{center}\rule{0.5\linewidth}{0.5pt}\end{center}

\section{\texorpdfstring{Common Error 3.3 -- Confusing \texttt{and} and
\texttt{or}}{Common Error 3.3 -- Confusing and and or}}\label{common-error-3.3-confusing-and-and-or}

Suppose a value is valid only when it lies between \texttt{0} and
\texttt{100}.

Correct:

\begin{Shaded}
\begin{Highlighting}[]
\ControlFlowTok{if}\NormalTok{ value }\OperatorTok{\textgreater{}=} \DecValTok{0} \KeywordTok{and}\NormalTok{ value }\OperatorTok{\textless{}=} \DecValTok{100}\NormalTok{ :}
    \BuiltInTok{print}\NormalTok{(}\StringTok{"Valid"}\NormalTok{)}
\end{Highlighting}
\end{Shaded}

Both conditions must hold.

To detect an invalid value:

\begin{Shaded}
\begin{Highlighting}[]
\ControlFlowTok{if}\NormalTok{ value }\OperatorTok{\textless{}} \DecValTok{0} \KeywordTok{or}\NormalTok{ value }\OperatorTok{\textgreater{}} \DecValTok{100}\NormalTok{ :}
    \BuiltInTok{print}\NormalTok{(}\StringTok{"Invalid"}\NormalTok{)}
\end{Highlighting}
\end{Shaded}

Only one of the invalid conditions needs to hold.

A helpful question is:

\begin{itemize}
\tightlist
\item
  Must \textbf{all} conditions be true? → use \texttt{and}.
\item
  Is \textbf{any one} sufficient? → use \texttt{or}.
\end{itemize}

\begin{center}\rule{0.5\linewidth}{0.5pt}\end{center}

\section{Programming Tip 3.4 --
Readability}\label{programming-tip-3.4-readability}

Avoid comparing Boolean variables explicitly with \texttt{True} or
\texttt{False}.

Less readable:

\begin{Shaded}
\begin{Highlighting}[]
\ControlFlowTok{if}\NormalTok{ frozen }\OperatorTok{==} \VariableTok{False}\NormalTok{ :}
    \BuiltInTok{print}\NormalTok{(}\StringTok{"Not frozen"}\NormalTok{)}
\end{Highlighting}
\end{Shaded}

Prefer:

\begin{Shaded}
\begin{Highlighting}[]
\ControlFlowTok{if} \KeywordTok{not}\NormalTok{ frozen :}
    \BuiltInTok{print}\NormalTok{(}\StringTok{"Not frozen"}\NormalTok{)}
\end{Highlighting}
\end{Shaded}

Likewise:

\begin{Shaded}
\begin{Highlighting}[]
\ControlFlowTok{if}\NormalTok{ valid :}
    \BuiltInTok{print}\NormalTok{(}\StringTok{"Input accepted"}\NormalTok{)}
\end{Highlighting}
\end{Shaded}

is usually clearer than:

\begin{Shaded}
\begin{Highlighting}[]
\ControlFlowTok{if}\NormalTok{ valid }\OperatorTok{==} \VariableTok{True}\NormalTok{ :}
    \BuiltInTok{print}\NormalTok{(}\StringTok{"Input accepted"}\NormalTok{)}
\end{Highlighting}
\end{Shaded}

Boolean variable names should describe a condition clearly:

\begin{Shaded}
\begin{Highlighting}[]
\NormalTok{valid}
\NormalTok{finished}
\NormalTok{found}
\NormalTok{isEmpty}
\NormalTok{hasError}
\end{Highlighting}
\end{Shaded}

\begin{center}\rule{0.5\linewidth}{0.5pt}\end{center}

\section{Special Topic 3.3 -- Chaining Relational
Operators}\label{special-topic-3.3-chaining-relational-operators}

Python allows:

\begin{Shaded}
\begin{Highlighting}[]
\DecValTok{0} \OperatorTok{\textless{}=}\NormalTok{ value }\OperatorTok{\textless{}=} \DecValTok{100}
\end{Highlighting}
\end{Shaded}

This is equivalent to:

\begin{Shaded}
\begin{Highlighting}[]
\NormalTok{value }\OperatorTok{\textgreater{}=} \DecValTok{0} \KeywordTok{and}\NormalTok{ value }\OperatorTok{\textless{}=} \DecValTok{100}
\end{Highlighting}
\end{Shaded}

Chained comparisons are concise and Pythonic, but beginners should
understand the equivalent Boolean expression explicitly.

\begin{center}\rule{0.5\linewidth}{0.5pt}\end{center}

\section{Special Topic 3.4 -- Short-Circuit
Evaluation}\label{special-topic-3.4-short-circuit-evaluation}

Python evaluates \texttt{and} and \texttt{or} from left to right and
stops as soon as the final result is known.

Example:

\begin{Shaded}
\begin{Highlighting}[]
\NormalTok{quantity }\OperatorTok{\textgreater{}} \DecValTok{0} \KeywordTok{and}\NormalTok{ price }\OperatorTok{/}\NormalTok{ quantity }\OperatorTok{\textless{}} \DecValTok{10}
\end{Highlighting}
\end{Shaded}

If \texttt{quantity\ \textgreater{}\ 0} is false, Python does not
evaluate the division. This prevents division by zero when
\texttt{quantity} is \texttt{0}.

For \texttt{or}:

\begin{Shaded}
\begin{Highlighting}[]
\NormalTok{conditionA }\KeywordTok{or}\NormalTok{ conditionB}
\end{Highlighting}
\end{Shaded}

if \texttt{conditionA} is already true, \texttt{conditionB} is not
evaluated.

This behavior is called \textbf{short-circuit evaluation}.

\begin{center}\rule{0.5\linewidth}{0.5pt}\end{center}

\section{Special Topic 3.5 -- De Morgan's
Laws}\label{special-topic-3.5-de-morgans-laws}

De Morgan's laws help rewrite negated compound conditions.

\begin{Shaded}
\begin{Highlighting}[]
\NormalTok{not (A and B)  ≡  (not A) or (not B)}

\NormalTok{not (A or B)   ≡  (not A) and (not B)}
\end{Highlighting}
\end{Shaded}

Example:

\begin{Shaded}
\begin{Highlighting}[]
\KeywordTok{not}\NormalTok{ (state }\OperatorTok{==} \StringTok{"AK"} \KeywordTok{or}\NormalTok{ state }\OperatorTok{==} \StringTok{"HI"}\NormalTok{)}
\end{Highlighting}
\end{Shaded}

is equivalent to:

\begin{Shaded}
\begin{Highlighting}[]
\NormalTok{state }\OperatorTok{!=} \StringTok{"AK"} \KeywordTok{and}\NormalTok{ state }\OperatorTok{!=} \StringTok{"HI"}
\end{Highlighting}
\end{Shaded}

Another example:

\begin{Shaded}
\begin{Highlighting}[]
\KeywordTok{not}\NormalTok{ (country }\OperatorTok{==} \StringTok{"USA"} \KeywordTok{and}\NormalTok{ state }\OperatorTok{!=} \StringTok{"AK"} \KeywordTok{and}\NormalTok{ state }\OperatorTok{!=} \StringTok{"HI"}\NormalTok{)}
\end{Highlighting}
\end{Shaded}

can be rewritten more directly as:

\begin{Shaded}
\begin{Highlighting}[]
\NormalTok{country }\OperatorTok{!=} \StringTok{"USA"} \KeywordTok{or}\NormalTok{ state }\OperatorTok{==} \StringTok{"AK"} \KeywordTok{or}\NormalTok{ state }\OperatorTok{==} \StringTok{"HI"}
\end{Highlighting}
\end{Shaded}

The goal is not merely shorter code; it is clearer logic.

\begin{center}\rule{0.5\linewidth}{0.5pt}\end{center}

\subsection{Self Check 3.7}\label{self-check-3.7}

\subsubsection{Question 1}\label{question-1-6}

How do you test whether both \texttt{x} and \texttt{y} are positive?

\textbf{Answer}

\begin{Shaded}
\begin{Highlighting}[]
\NormalTok{x }\OperatorTok{\textgreater{}} \DecValTok{0} \KeywordTok{and}\NormalTok{ y }\OperatorTok{\textgreater{}} \DecValTok{0}
\end{Highlighting}
\end{Shaded}

\subsubsection{Question 2}\label{question-2-6}

How do you test whether at least one of \texttt{x} and \texttt{y} is
zero?

\textbf{Answer}

\begin{Shaded}
\begin{Highlighting}[]
\NormalTok{x }\OperatorTok{==} \DecValTok{0} \KeywordTok{or}\NormalTok{ y }\OperatorTok{==} \DecValTok{0}
\end{Highlighting}
\end{Shaded}

\subsubsection{Question 3}\label{question-3-3}

What is the value of:

\begin{Shaded}
\begin{Highlighting}[]
\KeywordTok{not} \KeywordTok{not} \VariableTok{True}
\end{Highlighting}
\end{Shaded}

\textbf{Answer}

\texttt{True}.

\subsubsection{Question 4}\label{question-4-2}

Why can this be safe when \texttt{quantity\ ==\ 0}?

\begin{Shaded}
\begin{Highlighting}[]
\NormalTok{quantity }\OperatorTok{\textgreater{}} \DecValTok{0} \KeywordTok{and}\NormalTok{ price }\OperatorTok{/}\NormalTok{ quantity }\OperatorTok{\textless{}} \DecValTok{10}
\end{Highlighting}
\end{Shaded}

\textbf{Answer}

Because short-circuit evaluation stops after the first false condition,
so the division is not attempted.

\begin{center}\rule{0.5\linewidth}{0.5pt}\end{center}

\section{3.8 Analyzing Strings}\label{analyzing-strings}

Decision logic is often used to examine text.

Python provides operators and methods for checking whether a string
contains another string or has particular characteristics.

\begin{center}\rule{0.5\linewidth}{0.5pt}\end{center}

\subsection{\texorpdfstring{3.8.1 Membership with \texttt{in} and
\texttt{not\ in}}{3.8.1 Membership with in and not in}}\label{membership-with-in-and-not-in}

\begin{Shaded}
\begin{Highlighting}[]
\NormalTok{name }\OperatorTok{=} \StringTok{"John Wayne"}

\ControlFlowTok{if} \StringTok{"Way"} \KeywordTok{in}\NormalTok{ name :}
    \BuiltInTok{print}\NormalTok{(}\StringTok{"Substring found"}\NormalTok{)}
\end{Highlighting}
\end{Shaded}

The inverse test is:

\begin{Shaded}
\begin{Highlighting}[]
\ControlFlowTok{if} \StringTok{"{-}"} \KeywordTok{not} \KeywordTok{in}\NormalTok{ name :}
    \BuiltInTok{print}\NormalTok{(}\StringTok{"The name does not contain a hyphen"}\NormalTok{)}
\end{Highlighting}
\end{Shaded}

Membership tests are case-sensitive.

\begin{center}\rule{0.5\linewidth}{0.5pt}\end{center}

\subsection{3.8.2 Useful Substring
Methods}\label{useful-substring-methods}

\begin{longtable}[]{@{}
  >{\raggedright\arraybackslash}p{(\columnwidth - 2\tabcolsep) * \real{0.5000}}
  >{\raggedright\arraybackslash}p{(\columnwidth - 2\tabcolsep) * \real{0.5000}}@{}}
\toprule\noalign{}
\begin{minipage}[b]{\linewidth}\raggedright
Operation
\end{minipage} & \begin{minipage}[b]{\linewidth}\raggedright
Meaning
\end{minipage} \\
\midrule\noalign{}
\endhead
\bottomrule\noalign{}
\endlastfoot
\texttt{substring\ in\ s} & Does \texttt{s} contain
\texttt{substring}? \\
\texttt{s.count(substring)} & Number of non-overlapping occurrences \\
\texttt{s.startswith(substring)} & Does \texttt{s} begin with it? \\
\texttt{s.endswith(substring)} & Does \texttt{s} end with it? \\
\texttt{s.find(substring)} & Index of first occurrence, or \texttt{-1}
if not found \\
\end{longtable}

Example:

\begin{Shaded}
\begin{Highlighting}[]
\NormalTok{filename }\OperatorTok{=} \StringTok{"report.html"}

\ControlFlowTok{if}\NormalTok{ filename.endswith(}\StringTok{".html"}\NormalTok{) :}
    \BuiltInTok{print}\NormalTok{(}\StringTok{"HTML file"}\NormalTok{)}
\end{Highlighting}
\end{Shaded}

Example:

\begin{Shaded}
\begin{Highlighting}[]
\NormalTok{text }\OperatorTok{=} \StringTok{"banana"}
\BuiltInTok{print}\NormalTok{(text.count(}\StringTok{"an"}\NormalTok{))}
\BuiltInTok{print}\NormalTok{(text.find(}\StringTok{"na"}\NormalTok{))}
\end{Highlighting}
\end{Shaded}

\begin{center}\rule{0.5\linewidth}{0.5pt}\end{center}

\subsection{3.8.3 Testing String
Characteristics}\label{testing-string-characteristics}

Useful methods include:

\begin{longtable}[]{@{}
  >{\raggedright\arraybackslash}p{(\columnwidth - 2\tabcolsep) * \real{0.5000}}
  >{\raggedright\arraybackslash}p{(\columnwidth - 2\tabcolsep) * \real{0.5000}}@{}}
\toprule\noalign{}
\begin{minipage}[b]{\linewidth}\raggedright
Method
\end{minipage} & \begin{minipage}[b]{\linewidth}\raggedright
True when\ldots{}
\end{minipage} \\
\midrule\noalign{}
\endhead
\bottomrule\noalign{}
\endlastfoot
\texttt{s.isalnum()} & all characters are letters or digits and the
string is nonempty \\
\texttt{s.isalpha()} & all characters are letters and the string is
nonempty \\
\texttt{s.isdigit()} & all characters are digits and the string is
nonempty \\
\texttt{s.islower()} & there is at least one letter and all letters are
lowercase \\
\texttt{s.isupper()} & there is at least one letter and all letters are
uppercase \\
\texttt{s.isspace()} & all characters are whitespace and the string is
nonempty \\
\end{longtable}

Examples:

\begin{Shaded}
\begin{Highlighting}[]
\CommentTok{"1729"}\NormalTok{.isdigit()}
\end{Highlighting}
\end{Shaded}

returns \texttt{True}.

\begin{Shaded}
\begin{Highlighting}[]
\CommentTok{"{-}1729"}\NormalTok{.isdigit()}
\end{Highlighting}
\end{Shaded}

returns \texttt{False} because \texttt{-} is not a digit.

\begin{Shaded}
\begin{Highlighting}[]
\CommentTok{"John Smith"}\NormalTok{.isalpha()}
\end{Highlighting}
\end{Shaded}

returns \texttt{False} because the string contains a space.

\begin{center}\rule{0.5\linewidth}{0.5pt}\end{center}

\subsection{Example -- Exploring a
Substring}\label{example-exploring-a-substring}

\begin{Shaded}
\begin{Highlighting}[]
\NormalTok{theString }\OperatorTok{=} \BuiltInTok{input}\NormalTok{(}\StringTok{"Enter a string: "}\NormalTok{)}
\NormalTok{theSubstring }\OperatorTok{=} \BuiltInTok{input}\NormalTok{(}\StringTok{"Enter a substring: "}\NormalTok{)}

\ControlFlowTok{if}\NormalTok{ theSubstring }\KeywordTok{in}\NormalTok{ theString :}
    \BuiltInTok{print}\NormalTok{(}\StringTok{"The string contains the substring."}\NormalTok{)}
    \BuiltInTok{print}\NormalTok{(}\StringTok{"Count:"}\NormalTok{, theString.count(theSubstring))}
    \BuiltInTok{print}\NormalTok{(}\StringTok{"First position:"}\NormalTok{, theString.find(theSubstring))}

    \ControlFlowTok{if}\NormalTok{ theString.startswith(theSubstring) :}
        \BuiltInTok{print}\NormalTok{(}\StringTok{"It appears at the beginning."}\NormalTok{)}

    \ControlFlowTok{if}\NormalTok{ theString.endswith(theSubstring) :}
        \BuiltInTok{print}\NormalTok{(}\StringTok{"It appears at the end."}\NormalTok{)}
\ControlFlowTok{else}\NormalTok{ :}
    \BuiltInTok{print}\NormalTok{(}\StringTok{"The string does not contain the substring."}\NormalTok{)}
\end{Highlighting}
\end{Shaded}

\begin{center}\rule{0.5\linewidth}{0.5pt}\end{center}

\subsection{Self Check 3.8}\label{self-check-3.8}

\subsubsection{Question 1}\label{question-1-7}

How do you count blank spaces in \texttt{text}?

\textbf{Answer}

\begin{Shaded}
\begin{Highlighting}[]
\NormalTok{text.count(}\StringTok{" "}\NormalTok{)}
\end{Highlighting}
\end{Shaded}

\subsubsection{Question 2}\label{question-2-7}

How do you test whether a filename ends with \texttt{.jpg} or
\texttt{.jpeg}?

\textbf{Answer}

\begin{Shaded}
\begin{Highlighting}[]
\NormalTok{filename.endswith(}\StringTok{".jpg"}\NormalTok{) }\KeywordTok{or}\NormalTok{ filename.endswith(}\StringTok{".jpeg"}\NormalTok{)}
\end{Highlighting}
\end{Shaded}

For case-insensitive checking, first normalize the name with
\texttt{lower()}.

\subsubsection{Question 3}\label{question-3-4}

What does \texttt{find()} return when the substring is absent?

\textbf{Answer}

\texttt{-1}.

\begin{center}\rule{0.5\linewidth}{0.5pt}\end{center}

\section{3.9 Application: Input
Validation}\label{application-input-validation}

\textbf{Input validation} means checking user-supplied data before using
it.

A program should not assume that every input satisfies the required
rules.

Suppose an elevator accepts floors \texttt{1} through \texttt{20},
except \texttt{13}.

Invalid inputs include:

\begin{itemize}
\tightlist
\item
  \texttt{13};
\item
  \texttt{0} or a negative number;
\item
  a number greater than \texttt{20}.
\end{itemize}

A validated version is:

\begin{Shaded}
\begin{Highlighting}[]
\NormalTok{floor }\OperatorTok{=} \BuiltInTok{int}\NormalTok{(}\BuiltInTok{input}\NormalTok{(}\StringTok{"Floor: "}\NormalTok{))}

\ControlFlowTok{if}\NormalTok{ floor }\OperatorTok{==} \DecValTok{13}\NormalTok{ :}
    \BuiltInTok{print}\NormalTok{(}\StringTok{"Error: There is no thirteenth floor."}\NormalTok{)}
\ControlFlowTok{elif}\NormalTok{ floor }\OperatorTok{\textless{}=} \DecValTok{0} \KeywordTok{or}\NormalTok{ floor }\OperatorTok{\textgreater{}} \DecValTok{20}\NormalTok{ :}
    \BuiltInTok{print}\NormalTok{(}\StringTok{"Error: The floor must be between 1 and 20."}\NormalTok{)}
\ControlFlowTok{else}\NormalTok{ :}
\NormalTok{    actualFloor }\OperatorTok{=}\NormalTok{ floor}

    \ControlFlowTok{if}\NormalTok{ floor }\OperatorTok{\textgreater{}} \DecValTok{13}\NormalTok{ :}
\NormalTok{        actualFloor }\OperatorTok{=}\NormalTok{ floor }\OperatorTok{{-}} \DecValTok{1}

    \BuiltInTok{print}\NormalTok{(}\StringTok{"The elevator will travel to the actual floor"}\NormalTok{, actualFloor)}
\end{Highlighting}
\end{Shaded}

The key pattern is:

\begin{Shaded}
\begin{Highlighting}[]
\NormalTok{Read input}
\NormalTok{   ↓}
\NormalTok{Validate input}
\NormalTok{   ↓}
\NormalTok{Invalid? → report error}
\NormalTok{Valid?   → process data}
\end{Highlighting}
\end{Shaded}

\begin{center}\rule{0.5\linewidth}{0.5pt}\end{center}

\subsection{3.9.1 Range Validation}\label{range-validation}

Example:

\begin{Shaded}
\begin{Highlighting}[]
\ControlFlowTok{if}\NormalTok{ age }\OperatorTok{\textless{}} \DecValTok{0} \KeywordTok{or}\NormalTok{ age }\OperatorTok{\textgreater{}} \DecValTok{130}\NormalTok{ :}
    \BuiltInTok{print}\NormalTok{(}\StringTok{"Invalid age"}\NormalTok{)}
\ControlFlowTok{else}\NormalTok{ :}
    \CommentTok{\# Safe to process the value.}
\NormalTok{    ...}
\end{Highlighting}
\end{Shaded}

\begin{center}\rule{0.5\linewidth}{0.5pt}\end{center}

\subsection{3.9.2 Validating Choice
Inputs}\label{validating-choice-inputs}

Suppose the user must enter \texttt{s} or \texttt{m}.

One approach is:

\begin{Shaded}
\begin{Highlighting}[]
\NormalTok{status }\OperatorTok{=} \BuiltInTok{input}\NormalTok{(}\StringTok{"Enter s or m: "}\NormalTok{)}

\ControlFlowTok{if}\NormalTok{ status }\OperatorTok{==} \StringTok{"s"} \KeywordTok{or}\NormalTok{ status }\OperatorTok{==} \StringTok{"S"}\NormalTok{ :}
\NormalTok{    ...}
\ControlFlowTok{elif}\NormalTok{ status }\OperatorTok{==} \StringTok{"m"} \KeywordTok{or}\NormalTok{ status }\OperatorTok{==} \StringTok{"M"}\NormalTok{ :}
\NormalTok{    ...}
\ControlFlowTok{else}\NormalTok{ :}
    \BuiltInTok{print}\NormalTok{(}\StringTok{"Invalid status"}\NormalTok{)}
\end{Highlighting}
\end{Shaded}

A more scalable technique is to normalize first:

\begin{Shaded}
\begin{Highlighting}[]
\NormalTok{status }\OperatorTok{=} \BuiltInTok{input}\NormalTok{(}\StringTok{"Enter s or m: "}\NormalTok{)}
\NormalTok{status }\OperatorTok{=}\NormalTok{ status.lower()}

\ControlFlowTok{if}\NormalTok{ status }\OperatorTok{==} \StringTok{"s"}\NormalTok{ :}
\NormalTok{    ...}
\ControlFlowTok{elif}\NormalTok{ status }\OperatorTok{==} \StringTok{"m"}\NormalTok{ :}
\NormalTok{    ...}
\ControlFlowTok{else}\NormalTok{ :}
    \BuiltInTok{print}\NormalTok{(}\StringTok{"Invalid status"}\NormalTok{)}
\end{Highlighting}
\end{Shaded}

Similarly, multi-character codes can be normalized:

\begin{Shaded}
\begin{Highlighting}[]
\NormalTok{country }\OperatorTok{=} \BuiltInTok{input}\NormalTok{(}\StringTok{"Country code: "}\NormalTok{)}
\NormalTok{country }\OperatorTok{=}\NormalTok{ country.upper()}
\end{Highlighting}
\end{Shaded}

Then the program only needs to compare against one representation.

\begin{center}\rule{0.5\linewidth}{0.5pt}\end{center}

\subsection{3.9.3 A Limitation at This
Stage}\label{a-limitation-at-this-stage}

This code:

\begin{Shaded}
\begin{Highlighting}[]
\NormalTok{floor }\OperatorTok{=} \BuiltInTok{int}\NormalTok{(}\BuiltInTok{input}\NormalTok{(}\StringTok{"Floor: "}\NormalTok{))}
\end{Highlighting}
\end{Shaded}

will still raise an exception if the user types something such as:

\begin{Shaded}
\begin{Highlighting}[]
\NormalTok{five}
\end{Highlighting}
\end{Shaded}

Checking conversion failures requires exception handling, which is
introduced later in the textbook.

At this stage, distinguish between:

\begin{itemize}
\tightlist
\item
  \textbf{validating the value after conversion} --- for example,
  checking whether it lies in a range;
\item
  \textbf{handling a failed conversion} --- a later topic involving
  exceptions.
\end{itemize}

\begin{center}\rule{0.5\linewidth}{0.5pt}\end{center}

\section{Special Topic 3.6 -- Terminating a
Program}\label{special-topic-3.6-terminating-a-program}

For small text-based programs, invalid input can sometimes be handled by
terminating immediately.

\begin{Shaded}
\begin{Highlighting}[]
\ImportTok{from}\NormalTok{ sys }\ImportTok{import}\NormalTok{ exit}

\NormalTok{response }\OperatorTok{=} \BuiltInTok{input}\NormalTok{(}\StringTok{"Enter y or n: "}\NormalTok{)}

\ControlFlowTok{if} \KeywordTok{not}\NormalTok{ (response }\OperatorTok{==} \StringTok{"y"} \KeywordTok{or}\NormalTok{ response }\OperatorTok{==} \StringTok{"n"}\NormalTok{) :}
\NormalTok{    exit(}\StringTok{"Error: you must enter y or n."}\NormalTok{)}
\end{Highlighting}
\end{Shaded}

After validation succeeds, later code can assume the input satisfies the
required rule.

Use early termination carefully. In larger programs, functions and
exceptions often provide better structure, but those ideas are
introduced later.

\begin{center}\rule{0.5\linewidth}{0.5pt}\end{center}

\section{Special Topic 3.7 -- Interactive Graphical
Programs}\label{special-topic-3.7-interactive-graphical-programs}

Decision logic applies to graphical programs as well.

A graphical program may:

\begin{enumerate}
\def\labelenumi{\arabic{enumi}.}
\tightlist
\item
  read input;
\item
  validate the coordinates or dimensions;
\item
  draw only if the values are valid;
\item
  react to a mouse click and branch according to its location.
\end{enumerate}

The key lesson is that graphical input does not eliminate the need for
validation.

\begin{center}\rule{0.5\linewidth}{0.5pt}\end{center}

\section{Computing \& Society 3.2 -- Artificial
Intelligence}\label{computing-society-3.2-artificial-intelligence}

The chapter introduces artificial intelligence as a broader context for
decision-making software. Traditional programs follow explicitly
designed rules, while AI research aims to build systems capable of
behavior that appears more flexible or intelligent.

For this lesson, the connection is conceptual: even sophisticated
systems must ultimately represent information, evaluate alternatives,
and select actions. Later AI techniques may learn parts of those
decision rules from data, but understanding explicit control flow
remains foundational programming knowledge.

\begin{center}\rule{0.5\linewidth}{0.5pt}\end{center}

\section{Worked Example 3.2 -- Intersecting
Circles}\label{worked-example-3.2-intersecting-circles}

The chapter's graphical worked example determines how two circles relate
to one another.

Suppose the circles have centers:

\begin{Shaded}
\begin{Highlighting}[]
\NormalTok{(x0, y0)}
\NormalTok{(x1, y1)}
\end{Highlighting}
\end{Shaded}

and radii:

\begin{Shaded}
\begin{Highlighting}[]
\NormalTok{r0}
\NormalTok{r1}
\end{Highlighting}
\end{Shaded}

The distance between the centers is:

\begin{Shaded}
\begin{Highlighting}[]
\ImportTok{from}\NormalTok{ math }\ImportTok{import}\NormalTok{ sqrt}

\NormalTok{dist }\OperatorTok{=}\NormalTok{ sqrt((x1 }\OperatorTok{{-}}\NormalTok{ x0) }\OperatorTok{**} \DecValTok{2} \OperatorTok{+}\NormalTok{ (y1 }\OperatorTok{{-}}\NormalTok{ y0) }\OperatorTok{**} \DecValTok{2}\NormalTok{)}
\end{Highlighting}
\end{Shaded}

The program can classify the relationship using decisions:

\begin{Shaded}
\begin{Highlighting}[]
\ControlFlowTok{if}\NormalTok{ dist }\OperatorTok{\textgreater{}}\NormalTok{ r0 }\OperatorTok{+}\NormalTok{ r1 :}
\NormalTok{    message }\OperatorTok{=} \StringTok{"The circles are separate."}
\ControlFlowTok{elif}\NormalTok{ dist }\OperatorTok{\textless{}} \BuiltInTok{abs}\NormalTok{(r0 }\OperatorTok{{-}}\NormalTok{ r1) :}
\NormalTok{    message }\OperatorTok{=} \StringTok{"One circle is inside the other."}
\ControlFlowTok{elif}\NormalTok{ dist }\OperatorTok{==}\NormalTok{ r0 }\OperatorTok{+}\NormalTok{ r1 :}
\NormalTok{    message }\OperatorTok{=} \StringTok{"The circles touch externally."}
\ControlFlowTok{elif}\NormalTok{ dist }\OperatorTok{==} \DecValTok{0} \KeywordTok{and}\NormalTok{ r0 }\OperatorTok{==}\NormalTok{ r1 :}
\NormalTok{    message }\OperatorTok{=} \StringTok{"The circles coincide."}
\ControlFlowTok{else}\NormalTok{ :}
\NormalTok{    message }\OperatorTok{=} \StringTok{"The circles intersect at two points."}
\end{Highlighting}
\end{Shaded}

This example combines:

\begin{itemize}
\tightlist
\item
  input validation;
\item
  arithmetic;
\item
  relational operators;
\item
  Boolean operators;
\item
  multiple alternatives;
\item
  geometry;
\item
  graphical output.
\end{itemize}

For floating-point geometric computations, an approximate-equality test
may be preferable when values are not exact integers.

\begin{center}\rule{0.5\linewidth}{0.5pt}\end{center}

\section{TOOLBOX 3.2 -- Plotting Simple
Graphs}\label{toolbox-3.2-plotting-simple-graphs}

The chapter introduces \texttt{matplotlib.pyplot} as an optional example
of using a Python library to visualize data.

A simple bar chart:

\begin{Shaded}
\begin{Highlighting}[]
\ImportTok{from}\NormalTok{ matplotlib }\ImportTok{import}\NormalTok{ pyplot}

\NormalTok{pyplot.bar(}\DecValTok{1}\NormalTok{, }\FloatTok{1.1}\NormalTok{)}
\NormalTok{pyplot.bar(}\DecValTok{2}\NormalTok{, }\FloatTok{10.0}\NormalTok{)}
\NormalTok{pyplot.bar(}\DecValTok{3}\NormalTok{, }\FloatTok{25.4}\NormalTok{)}
\NormalTok{pyplot.bar(}\DecValTok{4}\NormalTok{, }\FloatTok{44.5}\NormalTok{)}
\NormalTok{pyplot.bar(}\DecValTok{5}\NormalTok{, }\FloatTok{61.0}\NormalTok{)}

\NormalTok{pyplot.xlabel(}\StringTok{"Month"}\NormalTok{)}
\NormalTok{pyplot.ylabel(}\StringTok{"Temperature"}\NormalTok{)}
\NormalTok{pyplot.show()}
\end{Highlighting}
\end{Shaded}

A simple line graph:

\begin{Shaded}
\begin{Highlighting}[]
\ImportTok{from}\NormalTok{ matplotlib }\ImportTok{import}\NormalTok{ pyplot}

\NormalTok{months }\OperatorTok{=}\NormalTok{ [}\DecValTok{1}\NormalTok{, }\DecValTok{2}\NormalTok{, }\DecValTok{3}\NormalTok{, }\DecValTok{4}\NormalTok{, }\DecValTok{5}\NormalTok{]}
\NormalTok{temperatures }\OperatorTok{=}\NormalTok{ [}\FloatTok{1.1}\NormalTok{, }\FloatTok{10.0}\NormalTok{, }\FloatTok{25.4}\NormalTok{, }\FloatTok{44.5}\NormalTok{, }\FloatTok{61.0}\NormalTok{]}

\NormalTok{pyplot.plot(months, temperatures)}
\NormalTok{pyplot.xlabel(}\StringTok{"Month"}\NormalTok{)}
\NormalTok{pyplot.ylabel(}\StringTok{"Temperature"}\NormalTok{)}
\NormalTok{pyplot.show()}
\end{Highlighting}
\end{Shaded}

The plotting toolbox is optional. The main chapter objective remains
decision-making and Boolean logic.

\begin{center}\rule{0.5\linewidth}{0.5pt}\end{center}

\section{Chapter 3 Summary}\label{chapter-3-summary}

\subsection{\texorpdfstring{\texttt{if}
Statements}{if Statements}}\label{if-statements}

\begin{itemize}
\tightlist
\item
  An \texttt{if} statement changes the path of execution according to a
  condition.
\item
  An \texttt{else} branch provides an alternative when the condition is
  false.
\item
  \texttt{elif} supports multiple mutually exclusive alternatives.
\item
  A compound statement has a header and an indented block.
\item
  Consistent indentation is part of Python syntax.
\item
  Avoid duplicated code across branches when common work can be moved
  outside.
\end{itemize}

\subsection{Relational Operators}\label{relational-operators-2}

\begin{Shaded}
\begin{Highlighting}[]
\OperatorTok{\textless{}}  \OperatorTok{\textless{}=}  \OperatorTok{\textgreater{}}  \OperatorTok{\textgreater{}=}  \OperatorTok{==}  \OperatorTok{!=}
\end{Highlighting}
\end{Shaded}

Remember:

\begin{Shaded}
\begin{Highlighting}[]
\NormalTok{=   assignment}
\NormalTok{==  equality test}
\end{Highlighting}
\end{Shaded}

Strings can also be compared, and their comparisons are case-sensitive.

Computed floating-point values should often be compared using a
tolerance rather than exact equality.

\subsection{Nested Branches}\label{nested-branches-1}

\begin{itemize}
\tightlist
\item
  An \texttt{if} statement can occur inside another branch.
\item
  Nesting is useful for decisions with several levels.
\item
  Indentation shows which inner decision belongs to which outer branch.
\item
  Hand-tracing helps verify nested logic.
\end{itemize}

\subsection{Multiple Alternatives}\label{multiple-alternatives-2}

Use:

\begin{Shaded}
\begin{Highlighting}[]
\ControlFlowTok{if}\NormalTok{ ... :}
\NormalTok{    ...}
\ControlFlowTok{elif}\NormalTok{ ... :}
\NormalTok{    ...}
\ControlFlowTok{else}\NormalTok{ :}
\NormalTok{    ...}
\end{Highlighting}
\end{Shaded}

Only the first true branch executes.

For overlapping thresholds, order tests carefully.

\subsection{Flowcharts}\label{flowcharts}

Flowcharts provide a visual representation of:

\begin{itemize}
\tightlist
\item
  tasks;
\item
  input/output;
\item
  decisions;
\item
  branches and joins.
\end{itemize}

Use structured branches and avoid arrows that jump into another branch.

\subsection{Test Cases}\label{test-cases}

Good tests include:

\begin{itemize}
\tightlist
\item
  at least one test for every branch;
\item
  boundary values;
\item
  nearby values around boundaries;
\item
  invalid values when validation is required.
\end{itemize}

Designing expected results before coding strengthens both the algorithm
and the test process.

\subsection{Boolean Logic}\label{boolean-logic}

\begin{Shaded}
\begin{Highlighting}[]
\VariableTok{True}
\VariableTok{False}
\KeywordTok{and}
\KeywordTok{or}
\KeywordTok{not}
\end{Highlighting}
\end{Shaded}

\begin{itemize}
\tightlist
\item
  \texttt{and} requires all component conditions to be true.
\item
  \texttt{or} requires at least one component condition to be true.
\item
  \texttt{not} reverses a Boolean value.
\item
  Python uses short-circuit evaluation.
\item
  De Morgan's laws help simplify negated compound conditions.
\end{itemize}

\subsection{String Analysis}\label{string-analysis}

Useful operations include:

\begin{Shaded}
\begin{Highlighting}[]
\NormalTok{substring }\KeywordTok{in}\NormalTok{ text}
\NormalTok{substring }\KeywordTok{not} \KeywordTok{in}\NormalTok{ text}
\NormalTok{text.count(...)}
\NormalTok{text.find(...)}
\NormalTok{text.startswith(...)}
\NormalTok{text.endswith(...)}
\NormalTok{text.isalpha()}
\NormalTok{text.isdigit()}
\NormalTok{text.isalnum()}
\NormalTok{text.islower()}
\NormalTok{text.isupper()}
\NormalTok{text.isspace()}
\end{Highlighting}
\end{Shaded}

\subsection{Input Validation}\label{input-validation}

Before processing user input:

\begin{Shaded}
\begin{Highlighting}[]
\NormalTok{Read}
\NormalTok{→ normalize when appropriate}
\NormalTok{→ validate}
\NormalTok{→ reject invalid values}
\NormalTok{→ process valid data}
\end{Highlighting}
\end{Shaded}

Range checks and valid-choice checks can be implemented with conditions.
Failed numeric conversions require exception handling, which is covered
later.

\begin{center}\rule{0.5\linewidth}{0.5pt}\end{center}

\section{Summary Quiz}\label{summary-quiz}

\subsection{Question 1}\label{question-1-8}

Which operator tests equality?

A. \texttt{=}\\
B. \texttt{==}\\
C. \texttt{!=}\\
D. \texttt{\textless{}=}

\textbf{Answer}

\textbf{B. \texttt{==}}

\subsection{Question 2}\label{question-2-8}

What happens when the condition of an \texttt{if/else} statement is
true?

A. Both branches execute.\\
B. Only the \texttt{else} branch executes.\\
C. The \texttt{if} branch executes and the \texttt{else} branch is
skipped.\\
D. The program always terminates.

\textbf{Answer}

\textbf{C}

\subsection{Question 3}\label{question-3-5}

What is the value of:

\begin{Shaded}
\begin{Highlighting}[]
\DecValTok{5} \OperatorTok{\textgreater{}=} \DecValTok{5}
\end{Highlighting}
\end{Shaded}

A. \texttt{True}\\
B. \texttt{False}\\
C. \texttt{5}\\
D. Error

\textbf{Answer}

\textbf{A. \texttt{True}}

\subsection{Question 4}\label{question-4-3}

Which condition tests whether \texttt{x} lies strictly between
\texttt{0} and \texttt{10}?

A. \texttt{x\ \textgreater{}\ 0\ or\ x\ \textless{}\ 10}\\
B. \texttt{x\ \textgreater{}\ 0\ and\ x\ \textless{}\ 10}\\
C. \texttt{x\ \textgreater{}=\ 0\ or\ x\ \textless{}=\ 10}\\
D. \texttt{not\ x}

\textbf{Answer}

\textbf{B}

\subsection{Question 5}\label{question-5}

What is the result of:

\begin{Shaded}
\begin{Highlighting}[]
\VariableTok{True} \KeywordTok{or} \VariableTok{False}
\end{Highlighting}
\end{Shaded}

\textbf{Answer}

\texttt{True}.

\subsection{Question 6}\label{question-6}

Why is this test risky after floating-point calculations?

\begin{Shaded}
\begin{Highlighting}[]
\ControlFlowTok{if}\NormalTok{ result }\OperatorTok{==}\NormalTok{ expected :}
\end{Highlighting}
\end{Shaded}

\textbf{Answer}

Because tiny roundoff errors may make two mathematically equal values
differ in their floating-point representations. A tolerance comparison
may be more appropriate.

\subsection{Question 7}\label{question-7}

Which structure is usually best for mutually exclusive grade categories?

A. Several independent \texttt{if} statements\\
B. \texttt{if/elif/else}\\
C. Only an \texttt{else} statement\\
D. No conditions

\textbf{Answer}

\textbf{B. \texttt{if/elif/else}}

\subsection{Question 8}\label{question-8}

What does this expression test?

\begin{Shaded}
\begin{Highlighting}[]
\CommentTok{".csv"} \KeywordTok{in}\NormalTok{ filename}
\end{Highlighting}
\end{Shaded}

\textbf{Answer}

Whether the exact substring \texttt{".csv"} occurs anywhere in
\texttt{filename}.

\subsection{Question 9}\label{question-9}

What does \texttt{text.find("abc")} return if \texttt{"abc"} is absent?

A. \texttt{False}\\
B. \texttt{None}\\
C. \texttt{-1}\\
D. Error

\textbf{Answer}

\textbf{C. \texttt{-1}}

\subsection{Question 10}\label{question-10}

What is a boundary test for:

\begin{Shaded}
\begin{Highlighting}[]
\ControlFlowTok{if}\NormalTok{ score }\OperatorTok{\textgreater{}=} \DecValTok{50}\NormalTok{ :}
\end{Highlighting}
\end{Shaded}

\textbf{Answer}

\texttt{50} is the key boundary. \texttt{49} and \texttt{51} are also
useful neighboring tests.

\subsection{Question 11}\label{question-11}

Why can the order of conditions matter in an \texttt{if/elif} chain?

\textbf{Answer}

Because Python stops at the first true branch. A broad condition placed
too early can prevent a more specific later condition from ever being
reached.

\subsection{Question 12}\label{question-12}

What is the main purpose of input validation?

A. Make the code longer.\\
B. Ensure user-supplied values satisfy the program's assumptions before
they are processed.\\
C. Convert all values to strings.\\
D. Remove all \texttt{if} statements.

\textbf{Answer}

\textbf{B}

\begin{center}\rule{0.5\linewidth}{0.5pt}\end{center}

\section{Review Exercises}\label{review-exercises}

\subsection{Exercise 1. Trace a Two-Way
Decision}\label{exercise-1.-trace-a-two-way-decision}

Determine the final value of \texttt{fee}:

\begin{Shaded}
\begin{Highlighting}[]
\NormalTok{age }\OperatorTok{=} \DecValTok{16}

\ControlFlowTok{if}\NormalTok{ age }\OperatorTok{\textless{}} \DecValTok{18}\NormalTok{ :}
\NormalTok{    fee }\OperatorTok{=} \DecValTok{5}
\ControlFlowTok{else}\NormalTok{ :}
\NormalTok{    fee }\OperatorTok{=} \DecValTok{10}
\end{Highlighting}
\end{Shaded}

\textbf{Answer}

\texttt{fee} is \texttt{5}.

\begin{center}\rule{0.5\linewidth}{0.5pt}\end{center}

\subsection{Exercise 2. Choose the Correct Relational
Operator}\label{exercise-2.-choose-the-correct-relational-operator}

Complete the condition so that the message is printed when
\texttt{temperature} is at most \texttt{0}.

\begin{Shaded}
\begin{Highlighting}[]
\ControlFlowTok{if}\NormalTok{ temperature \_\_\_ }\DecValTok{0}\NormalTok{ :}
    \BuiltInTok{print}\NormalTok{(}\StringTok{"Freezing or below"}\NormalTok{)}
\end{Highlighting}
\end{Shaded}

\textbf{Answer}

\begin{Shaded}
\begin{Highlighting}[]
\OperatorTok{\textless{}=}
\end{Highlighting}
\end{Shaded}

\begin{center}\rule{0.5\linewidth}{0.5pt}\end{center}

\subsection{Exercise 3. Rewrite Without
Duplication}\label{exercise-3.-rewrite-without-duplication}

Improve:

\begin{Shaded}
\begin{Highlighting}[]
\ControlFlowTok{if}\NormalTok{ quantity }\OperatorTok{\textgreater{}=} \DecValTok{10}\NormalTok{ :}
\NormalTok{    rate }\OperatorTok{=} \FloatTok{0.90}
\NormalTok{    total }\OperatorTok{=}\NormalTok{ rate }\OperatorTok{*}\NormalTok{ price }\OperatorTok{*}\NormalTok{ quantity}
\ControlFlowTok{else}\NormalTok{ :}
\NormalTok{    rate }\OperatorTok{=} \FloatTok{1.00}
\NormalTok{    total }\OperatorTok{=}\NormalTok{ rate }\OperatorTok{*}\NormalTok{ price }\OperatorTok{*}\NormalTok{ quantity}
\end{Highlighting}
\end{Shaded}

\textbf{Sample Answer}

\begin{Shaded}
\begin{Highlighting}[]
\ControlFlowTok{if}\NormalTok{ quantity }\OperatorTok{\textgreater{}=} \DecValTok{10}\NormalTok{ :}
\NormalTok{    rate }\OperatorTok{=} \FloatTok{0.90}
\ControlFlowTok{else}\NormalTok{ :}
\NormalTok{    rate }\OperatorTok{=} \FloatTok{1.00}

\NormalTok{total }\OperatorTok{=}\NormalTok{ rate }\OperatorTok{*}\NormalTok{ price }\OperatorTok{*}\NormalTok{ quantity}
\end{Highlighting}
\end{Shaded}

\begin{center}\rule{0.5\linewidth}{0.5pt}\end{center}

\subsection{Exercise 4. Test Cases}\label{exercise-4.-test-cases}

For:

\begin{Shaded}
\begin{Highlighting}[]
\ControlFlowTok{if}\NormalTok{ balance }\OperatorTok{\textless{}} \DecValTok{0}\NormalTok{ :}
\NormalTok{    status }\OperatorTok{=} \StringTok{"Overdrawn"}
\ControlFlowTok{else}\NormalTok{ :}
\NormalTok{    status }\OperatorTok{=} \StringTok{"OK"}
\end{Highlighting}
\end{Shaded}

propose three useful test values.

\textbf{Sample Answer}

\begin{Shaded}
\begin{Highlighting}[]
\NormalTok{{-}1   → Overdrawn}
\NormalTok{0    → OK       (boundary)}
\NormalTok{1    → OK}
\end{Highlighting}
\end{Shaded}

\begin{center}\rule{0.5\linewidth}{0.5pt}\end{center}

\subsection{Exercise 5. Boolean
Expressions}\label{exercise-5.-boolean-expressions}

Write conditions for:

\begin{enumerate}
\def\labelenumi{\arabic{enumi}.}
\tightlist
\item
  \texttt{x} and \texttt{y} are both positive.
\item
  At least one of \texttt{x} and \texttt{y} is negative.
\item
  \texttt{age} is between \texttt{18} and \texttt{65}, inclusive.
\item
  \texttt{answer} is neither \texttt{"y"} nor \texttt{"n"}.
\end{enumerate}

\textbf{Sample Answer}

\begin{Shaded}
\begin{Highlighting}[]
\NormalTok{x }\OperatorTok{\textgreater{}} \DecValTok{0} \KeywordTok{and}\NormalTok{ y }\OperatorTok{\textgreater{}} \DecValTok{0}
\NormalTok{x }\OperatorTok{\textless{}} \DecValTok{0} \KeywordTok{or}\NormalTok{ y }\OperatorTok{\textless{}} \DecValTok{0}
\NormalTok{age }\OperatorTok{\textgreater{}=} \DecValTok{18} \KeywordTok{and}\NormalTok{ age }\OperatorTok{\textless{}=} \DecValTok{65}
\NormalTok{answer }\OperatorTok{!=} \StringTok{"y"} \KeywordTok{and}\NormalTok{ answer }\OperatorTok{!=} \StringTok{"n"}
\end{Highlighting}
\end{Shaded}

\begin{center}\rule{0.5\linewidth}{0.5pt}\end{center}

\subsection{Exercise 6. Analyze a
String}\label{exercise-6.-analyze-a-string}

Given:

\begin{Shaded}
\begin{Highlighting}[]
\NormalTok{filename }\OperatorTok{=} \StringTok{"report.Final.PDF"}
\end{Highlighting}
\end{Shaded}

write a case-insensitive test for whether the filename ends in
\texttt{.pdf}.

\textbf{Answer}

\begin{Shaded}
\begin{Highlighting}[]
\NormalTok{filename.lower().endswith(}\StringTok{".pdf"}\NormalTok{)}
\end{Highlighting}
\end{Shaded}

\begin{center}\rule{0.5\linewidth}{0.5pt}\end{center}

\section{Practice Exercises}\label{practice-exercises}

\subsection{Exercise 1. Pass / Fail / Invalid
Score}\label{exercise-1.-pass-fail-invalid-score}

Read an exam score.

\begin{itemize}
\tightlist
\item
  If the score is below \texttt{0} or above \texttt{100}, print an
  error.
\item
  Otherwise, print \texttt{"Pass"} for scores of at least \texttt{50}
  and \texttt{"Fail"} for lower scores.
\end{itemize}

Draw a flowchart before writing Python.

\begin{center}\rule{0.5\linewidth}{0.5pt}\end{center}

\subsection{Exercise 2. Largest of Two
Numbers}\label{exercise-2.-largest-of-two-numbers}

Read two numbers and print:

\begin{itemize}
\tightlist
\item
  the larger value;
\item
  \texttt{"Equal"} if they are the same.
\end{itemize}

Use an \texttt{if/elif/else} statement.

\begin{center}\rule{0.5\linewidth}{0.5pt}\end{center}

\subsection{Exercise 3. Largest of Three
Numbers}\label{exercise-3.-largest-of-three-numbers}

Read three numbers and determine the largest.

First solve the task using nested decisions. Then consider whether the
structure can be made clearer with Boolean expressions or multiple
alternatives.

\begin{center}\rule{0.5\linewidth}{0.5pt}\end{center}

\subsection{Exercise 4. Shipping Rule}\label{exercise-4.-shipping-rule}

Read:

\begin{itemize}
\tightlist
\item
  destination country;
\item
  state or province.
\end{itemize}

Normalize both inputs with \texttt{upper()}.

Use the following rules:

\begin{itemize}
\tightlist
\item
  continental USA: \texttt{\$5};
\item
  Alaska or Hawaii: \texttt{\$10};
\item
  outside the USA: \texttt{\$15}.
\end{itemize}

Print the shipping charge with two decimal places.

\begin{center}\rule{0.5\linewidth}{0.5pt}\end{center}

\subsection{Exercise 5. Grade
Classification}\label{exercise-5.-grade-classification}

Read a score from \texttt{0} to \texttt{100} and assign:

\begin{Shaded}
\begin{Highlighting}[]
\NormalTok{90–100 → A}
\NormalTok{80–89  → B}
\NormalTok{70–79  → C}
\NormalTok{60–69  → D}
\NormalTok{0–59   → F}
\end{Highlighting}
\end{Shaded}

Reject invalid scores before classifying them.

Create test cases for:

\begin{Shaded}
\begin{Highlighting}[]
\NormalTok{{-}1, 0, 59, 60, 69, 70, 79, 80, 89, 90, 100, 101}
\end{Highlighting}
\end{Shaded}

\begin{center}\rule{0.5\linewidth}{0.5pt}\end{center}

\subsection{Exercise 6. Username
Validation}\label{exercise-6.-username-validation}

Read a username and check that:

\begin{itemize}
\tightlist
\item
  it is not empty;
\item
  it contains only letters and digits;
\item
  it contains at least one character.
\end{itemize}

Use string-analysis methods from Section 3.8.

\begin{center}\rule{0.5\linewidth}{0.5pt}\end{center}

\subsection{Exercise 7. File-Type
Classifier}\label{exercise-7.-file-type-classifier}

Read a filename and classify it as:

\begin{itemize}
\tightlist
\item
  Python source: \texttt{.py};
\item
  notebook: \texttt{.ipynb};
\item
  CSV data: \texttt{.csv};
\item
  text file: \texttt{.txt};
\item
  other.
\end{itemize}

The check should be case-insensitive.

\begin{center}\rule{0.5\linewidth}{0.5pt}\end{center}

\subsection{Exercise 8. Approximate
Equality}\label{exercise-8.-approximate-equality}

Read two floating-point numbers and determine whether they differ by
less than \texttt{0.001}.

Do not use exact equality for the approximate comparison.

\begin{center}\rule{0.5\linewidth}{0.5pt}\end{center}

\subsection{Exercise 9. Triangle Type from Side
Lengths}\label{exercise-9.-triangle-type-from-side-lengths}

Read three positive side lengths.

First validate that they can form a triangle:

\begin{Shaded}
\begin{Highlighting}[]
\NormalTok{a + b \textgreater{} c}
\NormalTok{and}
\NormalTok{b + c \textgreater{} a}
\NormalTok{and}
\NormalTok{c + a \textgreater{} b}
\end{Highlighting}
\end{Shaded}

If valid, classify the triangle as:

\begin{itemize}
\tightlist
\item
  equilateral;
\item
  isosceles;
\item
  scalene.
\end{itemize}

Prepare branch and boundary tests before coding.

\begin{center}\rule{0.5\linewidth}{0.5pt}\end{center}

\subsection{Exercise 10. Simple Login-Code
Check}\label{exercise-10.-simple-login-code-check}

Read a short access code.

Validation rules:

\begin{itemize}
\tightlist
\item
  exactly six characters;
\item
  only letters and digits;
\item
  must contain the substring \texttt{"NEU"} ignoring case.
\end{itemize}

Print a specific message for valid or invalid input.

\begin{center}\rule{0.5\linewidth}{0.5pt}\end{center}

\section{Open Exercises}\label{open-exercises}

\subsection{Exercise 1}\label{exercise-1}

Give three real-world examples where a decision has exactly two
outcomes. Express each rule in pseudocode.

\subsection{Exercise 2}\label{exercise-2}

Give a real-world example that requires at least four mutually exclusive
alternatives. Explain why an \texttt{if/elif/else} structure is
appropriate.

\subsection{Exercise 3}\label{exercise-3}

Choose one condition involving a range and identify:

\begin{itemize}
\tightlist
\item
  a typical valid value;
\item
  a lower-boundary value;
\item
  an upper-boundary value;
\item
  a value just below the range;
\item
  a value just above the range.
\end{itemize}

\subsection{Exercise 4}\label{exercise-4}

Create an example in which independent \texttt{if} statements are
correct but an \texttt{if/elif} chain would be wrong. Explain why more
than one action may need to occur.

\subsection{Exercise 5}\label{exercise-5}

Write a complicated Boolean condition involving \texttt{not},
\texttt{and}, and \texttt{or}. Then simplify it using De Morgan's laws.

\begin{center}\rule{0.5\linewidth}{0.5pt}\end{center}

\section{Key Terms}\label{key-terms}

\begin{longtable}[]{@{}
  >{\raggedright\arraybackslash}p{(\columnwidth - 2\tabcolsep) * \real{0.5000}}
  >{\raggedright\arraybackslash}p{(\columnwidth - 2\tabcolsep) * \real{0.5000}}@{}}
\toprule\noalign{}
\begin{minipage}[b]{\linewidth}\raggedright
Term
\end{minipage} & \begin{minipage}[b]{\linewidth}\raggedright
Meaning
\end{minipage} \\
\midrule\noalign{}
\endhead
\bottomrule\noalign{}
\endlastfoot
Decision & Choosing an action according to a condition \\
Condition & Expression evaluated as true or false \\
\texttt{if} statement & Executes a block when a condition is true \\
\texttt{else} branch & Alternative block when an \texttt{if} condition
is false \\
\texttt{elif} & Adds another condition to a multi-way decision \\
Compound statement & Statement with a header and one or more nested
statements \\
Statement block & Group of statements at the same indentation level \\
Indentation & Leading spaces that define block structure in Python \\
Relational operator & Operator that compares two values \\
Equality test & Comparison using \texttt{==} \\
Boundary value & Value at the edge of a decision range \\
Floating-point tolerance & Small acceptable difference used for
approximate comparison \\
Nested branch & Decision statement inside another branch \\
Hand-tracing & Manually simulating program execution \\
Multiple alternatives & Decision with more than two possible branches \\
Flowchart & Diagram showing tasks and control flow \\
Branch coverage & Testing each possible branch of a decision \\
Test case & Input plus expected behavior used to verify a program \\
Boolean & Data type with values \texttt{True} and \texttt{False} \\
Flag & Boolean variable representing a condition or state \\
Boolean operator & \texttt{and}, \texttt{or}, or \texttt{not} \\
Short-circuit evaluation & Stopping Boolean evaluation once the result
is known \\
De Morgan's laws & Rules for transforming negated
\texttt{and}/\texttt{or} expressions \\
Substring & String contained within another string \\
Membership test & Test using \texttt{in} or \texttt{not\ in} \\
Input validation & Checking that input satisfies required constraints \\
Normalization & Converting input to a consistent form before
comparison \\
\end{longtable}

\begin{center}\rule{0.5\linewidth}{0.5pt}\end{center}

\section{Study Guidance}\label{study-guidance}

After Lesson 3, students should no longer think of a program as only a
fixed list of calculations. A program can now \textbf{choose its path}.

A useful workflow for decision problems is:

\begin{Shaded}
\begin{Highlighting}[]
\NormalTok{Understand the rule}
\NormalTok{→ identify the possible outcomes}
\NormalTok{→ formulate the condition(s)}
\NormalTok{→ check boundary values}
\NormalTok{→ draw a flowchart when useful}
\NormalTok{→ write pseudocode}
\NormalTok{→ choose if / if{-}else / if{-}elif{-}else / nested if}
\NormalTok{→ remove duplicated work}
\NormalTok{→ prepare test cases for every branch}
\NormalTok{→ implement in Python}
\NormalTok{→ run ordinary, boundary, and invalid{-}input tests}
\NormalTok{→ trace unexpected behavior}
\end{Highlighting}
\end{Shaded}

A useful practice cycle is:

\begin{quote}
\textbf{Predict → Trace → Run → Compare → Explain → Modify → Test
boundaries → Practice}
\end{quote}

The most important habit from this chapter is to treat every decision as
something that must be \textbf{designed and tested}, not merely typed.

\end{document}
